\documentclass[a4paper,11pt]{article}
\pdfoutput=1 % if your are submitting a pdflatex (i.e. if you have
             % images in pdf, png or jpg format)
\usepackage{jheppub} % for details on the use of the package, please
                     % see the JHEP-author-manual

%\usepackage[T1]{fontenc} % if needed
%\usepackage[utf8]{inputenc}

%\begin{figure}
%     \centering
%     \includegraphics[width=\textwidth{.png}
%     \caption{}
%     \label{fig:}
%\end{figure}

\usepackage{amssymb}
\usepackage{amsmath}
\usepackage{setspace}
\usepackage{amsfonts}
\usepackage{xcolor}
\usepackage{color}
\usepackage{braket}
\usepackage{bbm}
\usepackage{mathtools}
\usepackage{caption}
\usepackage{subcaption}
\usepackage{enumerate}
\usepackage{bm}
\usepackage[final]{changes}
%\setremarkmarkup{(#2)}
\definecolor{dolphingreen}{rgb}{0,0.74,0.63}
\definecolor{dolphinorange}{rgb}{0.98,0.53,0.098}
\definecolor{purple}{rgb}{0.4,.2,0.7}
\interfootnotelinepenalty=10000
%\usepackage[colorinlistoftodos]{todonotes}
\usepackage{mathrsfs}

% SHOW EQUATION LABELS
%\usepackage[inline]{showlabels}

% COMMENT IN AND OUT TO TURN OFF COMMENTS
%\newcommand{\comment}[1]{\textcolor{red}{#1}}
%\newcommand{\comment}[1]{\ignorespaces}
%\newcommand{\cts}[1]{\textcolor{magenta}{#1}}
\newcommand{\cts}[1]{\ignorespaces}

\usepackage{graphicx}
\graphicspath{{./Figs/}}
\usepackage{cleveref}

\DeclareMathOperator{\Tr}{Tr}
\newcommand{\la}{\langle}
\newcommand{\ra}{\rangle}
\newcommand{\tr}{\textrm{Tr}\,}
\newcommand{\norm}[1]{\left\|#1\right\|}
\newcommand{\fn}[1]{\footnote{#1}}
\newcommand{\qu}[1]{\textcolor{red}{(#1)}}
\newcommand{\st}[1]{\subsubsection*{\textcolor{blue}{#1}}}
\def\bne{\begin{equation}}
\def\ene{\end{equation}}
\newcommand{\del}{\partial}
\newcommand{\overbar}[1]{\mkern 1.5mu\overline{\mkern-1.5mu#1\mkern-1.5mu}\mkern 1.5mu}

\newsavebox{\imagebox}

\title{Island mirages}

%\author[a]{Author}
\author[a]{Andrew Rolph}

%\affiliation[a]{Affiliation}
\affiliation[a]{Institute for Theoretical Physics, University of Amsterdam, 1090 GL Amsterdam, The Netherlands}

%\emailAdd{email@gmail.com}
\emailAdd{andrew.d.rolph@googlemail.com}

\abstract{
We point out a loophole in the proof that the Island Finder conditions~\cite{Bousso2021} are sufficient for the existence of islands. 
We explore examples which satisfy the conditions, but have no islands, which we call island mirages.
We also describe and give resolutions to two new puzzles associated with islands, including a direct tension with the quantum Bousso bound.  
}



\begin{document}

%\iffalse
%\begin{Large} To do \end{Large}
%\begin{itemize}
%\item
%\end{itemize}
%\clearpage
%\newpage
%\fi 


\maketitle
\flushbottom

\section{Introduction}
%\cts{
%Outline of intro:
%\begin{itemize}
%\item BHIP, Page curve, holographic entanglement entropy, islands, necessary and sufficient conditions, QFC
%\end{itemize}
%Summary of what I do:
%\begin{itemize}
%\item In this paper we ...
%\item Loopholes in island finder
%\item Tension and resolution with quantum Bousso bound
%\item Tension and resolution of vacuum state violating the Bekenstein bound.
%\end{itemize}
%}
The black hole information problem (BHIP) is a tension between quantum mechanical unitarity and the insensitivity of Hawking radiation to initial conditions, and has been a persistent and impelling puzzle in quantum gravity since its conception~\cite{Hawking1975,Hawking1976}. The Page curve, the von Neumann entropy of emitted radiation as a function of time, is one sharp diagnostic of unitarity in black hole evaporation~\cite{Page1993}. Deriving a Page curve consistent with unitarity, one that ends at zero von Neumann entropy for a black hole formed from collapse of matter in a pure state, is an important step in resolving the BHIP, though it is not the whole story.  

Much of the recent progress in deriving a Page curve consistent with unitarity has been in developing the techniques to calculate the entropy of Hawking radiation, or more generally of quantum systems coupled to gravity. This was inspired in part by the Ryu-Takayanagi formula for calculating von Neumann entropy in holographic CFTs~\cite{Ryu_2006,Freedman2017}, and its generalisations~\cite{Hubeny2007,Engelhardt2014,Harper2018,Rolph2021}, and culminated in the island formula~\cite{Almheiri2019b} and quantum maximin prescription~\cite{Akers2019}. The island formula is a prescription for computing fine-grained entropy in quantum systems coupled to gravity in terms of semiclassical entropy. See~\cite{Almheiri2021} for a review of islands, including an explanation of what is meant by fine-grained, coarse-grained, renormalised and semiclassical entropy in this context. The island prescription is
\bne \label{eq:IslandFormula} S(R) = \min \text{ext}_I
 \left ( \frac{\text{Area}(\del I)}{4G_N} + S_{ren.}(R \cup I) \right ) + S_{c.t.} (R) \ene
where $S_{ren.}$ is the renormalised semiclassical von Neumann entropy, and $S_{c.t.}$ the entropy counterterm that isolates the UV divergences associated to the region's boundary. 
%\cts{Do I need to include counterterm?} 
The prescription instructs us to extremise a functional over all possible islands $I$, including the trivial empty island $I = \varnothing$, then pick the minimal extremum.

The main technical obstacle to explicitly finding islands is evaluating $S_{ren.} (R\cup I)$ on the set of all possible islands. The cases where it is possible, such as in lower dimensional and doubly holographic models~\cite{Almheiri2019b,Almheiri2019,Almheiri2019a,Almheiri2020}, are sparse. This motivates finding necessary and/or sufficient conditions for the existence of islands. Sets of necessary and/or sufficient conditions for the existence of non-trivial islands may be easier to evaluate than the island formula~\eqref{eq:IslandFormula} itself. A set of necessary conditions were proposed in~\cite{Hartman2020}, and a set of sufficient conditions in the Island Finder paper~\cite{Bousso2021}. The literature which relies on these sets of conditions is growing, so it is important to know if and when they can be relied upon.

The two conditions given in~\cite{Bousso2021} for a region $R$ to have a non-empty island are in terms of an `island detector'\footnote{Credit to Ahmed Almheiri for the name island detector.} region $I'$ and require
%\bne \begin{split}
% 1.& \quad  S_{gen} (I' \cup R) < S_{gen} (R) \\
% \text{and either }  2.a.& \pm \Theta_\pm (I' \cup R) \geq 0 \\
% \text{ or } 2.b.&  \pm \Theta_\pm (I' \cup R) \leq 0
%\end{split} \ene
\bne \label{eq:1aa} \qquad 1. \quad  S_{gen} (I' \cup R) < S_{gen} (R) \ene
and either
\bne  2.a. \label{eq:2aa}\quad \pm \Theta_\pm (I' \cup R) \geq 0 \ene
or
\bne 2.b. \label{eq:2bb}\quad  \pm \Theta_\pm (I' \cup R) \leq 0. \ene
where $\Theta_\pm$ is the quantum expansion in the future-directed outward $(+)$ and inward $(-)$ null directions normal to $\del I'$\footnote{See~\cite{Bousso2015} for a precise definition of quantum expansion.}.

In this paper we first point out a loophole in the proof that \cref{eq:1aa,eq:2aa,eq:2bb} are sufficient conditions for $R$ to have a non-empty island.
The proof given in~\cite{Bousso2021} is by contradiction: it is assumed that (1) there is an island detector region $I'$ satisfying conditions \eqref{eq:1aa} to \eqref{eq:2bb}, and (2) that the island $I$ is empty. Then it is shown that these two assumptions lead to a contradiction. The loophole is the implicit assumption that at least one of the time slices on which the empty island is the quantum maximin surface intersects the domain of dependence of $I'$\footnote{Assuming $I' \cup R$ is quantum normal, i.e.~\eqref{eq:2aa}. We will consider the quantum anti-normal case separately.}. If this assumption is \replaced{violated}{invalid} then the representative $\tilde{I}'$ - the intersection of $D(I')$ with the time slice - is itself empty and an intermediate step in the proof is invalid.  

We preempt and address counterarguments to this loophole which assert that the implicit assumption mentioned above is always valid. One counterargument is based on the claim that the past and future tips of $D(I')$ are truncated by singularities, and the basis of this claim are two puzzling contradictions that arise when considering small representatives $\tilde{I}'$ near either tip of $D(I')$.

The first puzzle is a direct tension with the quantum Bousso bound. Assuming that either of $\del D^\pm (I')$ is a quantum lightsheet, then the quantum Bousso bound implies that $S_{gen}(\tilde{I}' \cup R)$ must decrease as the representative $\tilde{I}'$ is deformed towards the tip of the $D^\pm (I')$. A contradiction, $S_{gen} (R) < S_{gen} (R)$, is found when the tip is reached and $\tilde{I}' = \varnothing$. We discuss several resolutions to this puzzle.
\deleted{(1) truncation of $D(I')$ by singularities, (2) $\del D^\pm (I')$ failing to be a quantum lightsheet due to caustics, and (3) to reach $\tilde {I}' = \varnothing$ one has to pass through UV regimes where the semiclassical relations like the quantum Bousso bound are invalid.}
\added{The first two are an evasion of the tension with the quantum Bousso bound: (1) there is no $\tilde{I}' = \varnothing$ because of truncation of $D(I')$ by curvature singularities, or (2) $\del D^\pm (I')$ fails to be a quantum lightsheet. The third resolution is that (3) the quantum Bousso bound \textit{is} violated, but that this is not unexpected because to reach $\tilde {I}' = \varnothing$ one has to pass outside the regime of validity of the semiclassical approximation.}
There is more than one resolution, and how the tension is resolved depends on the set-up.

The second puzzle is that the existence of island detectors naively and incorrectly implies that a ball-shaped region of the Minkowski vacuum violates the Bekenstein area bound. The authors of~\cite{Hartman2020} showed, with some mild assumptions, that any region $\mathcal{I}$ that satisfies $S_{gen.} (\mathcal{I} \cup R) < S_{gen.} (R)$, which includes island detector representatives, also violates the Bekenstein area bound. The second puzzle arises when applying this result to a representative that is smaller than the local curvature and temperature scales. If such a representative $\tilde{I}'$ exists for a given island detector $I'$, then the state on it is approximately that of the vacuum in flat space, which should not violate the Bekenstein area bound. 

\replaced{The puzzles are of broader importance than may initially appear, because they also need to be resolved for islands. As far as the puzzles are concerned islands are simply special cases of island detectors in the sense that they satisfy \eqref{eq:1aa} and saturate both conditions~\eqref{eq:2aa} and~\eqref{eq:2bb} by their definitional extremality.}
{The same puzzles need to be resolved for islands, even though they first arose when considering counterarguments to the loophole, because as far as the puzzles are concerned islands are simply special cases of island detectors in the sense that they satisfy \eqref{eq:1aa} and saturate both conditions~\eqref{eq:2aa} and~\eqref{eq:2bb} by their definitional extremality.}

We discuss in depth the resolution of these two puzzles for islands in asymptotically flat evaporating black holes. The past domain of dependence of the island has no curvature singularities to prevent one from considering arbitrarily small representatives\added{, so the quantum Bousso bound \textit{is} apparently violated}. A key part of resolving the Bekenstein bound puzzle is the sub-Planckian radial separation of the boundary of the island and the past lightcone of the radiation region\added{, which is an obstacle to specifying `nice' time slices with trace extrinsic curvature $|K| \ll l_p^{-1}$}.

%For an asymptotically flat evaporating black hole we argue that $D^+ (I)$ is truncated by a singularity, while for $D^- (I)$ 
%%it is resolved by the blueshift of the interior Hawking modes as they propagate towards the past.  
%Such a problematic representative also exists for an asymptotically flat evaporating black hole towards the past tip of $D^- (I)$, and we argue that the puzzle is resolved by the sub-Planckian radial separation of the boundary of the island and the past lightcone of the radiation region.

\textbf{Outline}: In Sec.~\ref{sec:IslandFinder} we review Island Finder and explain the loophole in the proof. In Sec.~\ref{sec:Puzzles} we discuss counterarguments to the loophole, which naturally leads us to the two puzzles, to which we propose resolutions. In Sec.~\ref{sec:Discussion} we give some open questions unresolved by this work and ideas for future research. In Appendix~\ref{sec:Mirages} we outline two set-ups which may have explicit examples of island mirages. 

\section{Island Finder}\label{sec:IslandFinder}
%\comment{This section I think is 80\% there, though that last 20\%  is important as a valid counterargument to my loophole would be deadly. 
%\begin{itemize}
%\item I would like to understand the blueshift effects discussed in section~\ref{sec:Quantum_Bousso_bound} better. \item There's also work to do in resolving the Bekenstein bound paradox of section~\ref{sec:BekBoundViol}, though it is not critical as it only needs to be resolved in cases where the quantum Bousso bound is applicable to the whole of $\del D^{\pm} (I')$, as I discuss.
%\end{itemize}}

In this section we discuss the Island Finder conditions and reproduce the proof given that they are sufficient conditions for the existence of a non-empty island~\cite{Bousso2021}. We point out loopholes in the proof. 

\subsection{Conditions}

\begin{figure}
     \centering
     \includegraphics[width=0.8\textwidth]{Figs/IslandProof.png}
     \caption{Two (asymptotically flat) quantum systems, one with dynamical gravity and one without. $R$ is a region in the non-gravitating system, and $I'$ a region in the gravitating system. $\Sigma$ is a time slice, and $\tilde{I}'$ the representative of $I'$ on that time slice as defined in~\eqref{eq:tildeI'def}. If $I'$ satisfies conditions \eqref{eq:cond1a} and \eqref{eq:cond2} then it is an island detector for $R$. \deleted{A loophole in the proof given in~\cite{Bousso2021} is the implicit assumption that $\Sigma$ intersects the domain of dependence $D(I')$.}}
     \label{fig:IslandProof}
\end{figure}

%Here, to make this paper more self-contained, we reproduce the mathematical argument underlying the Island Finder proof. 

Consider two disjoint achronal regions, $I'$ and $R$, in a spacetime $M$ with $R$ in a weakly gravitating subregion. 
%Take an achronal subregion $I'$ of a quantum system coupled to dynamical gravity, and an achronal subregion $R$ of a non-gravitational system. 
%Consider two quantum systems, one on a fixed background geometry and one coupled to dynamical gravity. 
%\cts{Need to comment on the coupled vs decoupled systems. And a precise way to say that $R$ is in a non or weakly gravitating region.} 
%Specify one achronal subregion in each system: $I'$ in the system with dynamical gravity and $R$ in the system without.  
The claim of \cite{Bousso2021} is that if two conditions are satisfied by the generalised entropy of $I' \cup R$ then this is sufficient to guarantee that $R$ has a non-empty island somewhere in $M$. 
%\footnote{Since the conditions only establish existence of the island, and do not say where the island is, a more appropriate name than Island Finder may have been Island Detector\cite{}}.

We will call an $I'$ that satisfies the two conditions an island detector region for $R$. If an island detector region exists and $R$ does have a non-empty island $I$, then the conditions are working as intended. If an island detector region exists and the island $I$ for $R$ \textit{is} empty, then we have a faulty island detector, and we call such an $I'$ an island mirage.

The first of the two conditions is
\bne \label{eq:cond1a} (1)\quad  S_{gen} (I' \cup R) < S_{gen}(R) \ene
%We will refer to this as the hyperentropic condition, since it implies that $I'$ is Bekenstein area bound violating~\cite{}, though this is strictly speaking a misnomer as \eqref{eq:cond1a} is a stronger condition than Bekenstein bound violation.
There is an equivalent form of this inequality, 
%\cts{Add references to finite $S_{gen}$, finiteness of MI for disconnected regions, and whoever wrote the condition this way first.}
\bne \label{eq:cond1b}
I(I',R) > S_{gen} (I') 
\ene 
which has the benefit of both sides being UV-finite even if we decouple gravity, $G_N \to 0$, in the region that contains $R$~\cite{Casini2007, Bombelli1986, Hartman2020}. One quick comment is $S_{gen}(I')$ scales as the surface area of $I'$ in Planck units, so there is a manifest need for a large degree of correlation between $I'$ and $R$ in order to satisfy the condition. 

The second of the two sufficient conditions is that $I'\cup R$ be either quantum anti-normal or quantum normal. Quantum normality is the property that the generalised entropy of a spatial region increases under deformations of the region's boundary in outward null directions. This is a generalisation of the classical notion of a normal spatial region whose surface area increases under outward deformations. The definition of quantum normality is 
\bne \label{eq:cond2}
(2) \quad \pm \Theta_\pm (I' \cup R) \geq 0
\ene
where $\Theta_\pm$ is the quantum expansion in the outward/inward future null directions~\cite{Bousso2015}. Quantum anti-normality is the opposite of quantum normality in the sense that the generalised entropy \textit{decreases} under outward null deformations. Any subregion of a time reflection symmetric slice is automatically quantum normal or anti-quantum normal, so any $I' \cup R$ on such a slice which satisfies condition \eqref{eq:cond1a} satisfies both of the Island Finder conditions. Also, if $I' \cup R$ is quantum normal then so is $I'$, which means the detector region $I'$ satisfies one of the necessary conditions to itself be an island~\cite{Almheiri2019a, Hartman2020}.  Note that \eqref{eq:cond2} is a property that is required to hold locally at every point on the boundary of $I' \cup R$, and that it depends on the non-local entanglement entropy.   



\subsection{QFC and quantum maximin}

In the proof the quantum focussing conjecture (QFC) and quantum maximin prescription are assumed~\cite{Bousso2015, Akers2019}. We assume that the spacetime is stably causal, because that is the causality condition assumed for the quantum maximin prescription~\cite{Akers2019}, and because we do not want to exclude asymptotically AdS spacetimes\footnote{Classical maximin was originally introduced as a reformulation of the HRT holographic entanglement entropy prescription which applies to asymptotically locally AdS spacetimes~\cite{Hubeny2007, Wall2012}. Since AdS is not globally hyperbolic it has no Cauchy slices, so the notion of Cauchy slices had to be enlarged to AdS-Cauchy slices which are geodesically complete achronal surfaces~\cite{Wall2012, Engelhardt2014}}. Stably causal spacetimes have a global time function, and the function's level sets are constant time slices~\cite{Wald:1984rg}. 

Quantum maximin is a prescription for calculating entropy of subregions in non-holographic systems as well as holographic ones~\cite{Akers2019}. A relevant class of models are AdS black holes evaporating into a non-gravitating bath, where in the non-gravitational description we couple a holographic and non-holographic system together~\cite{Almheiri2019a}. If $R$ is a subregion of the non-holographic system, then the quantum maximin prescription for calculating its exact von Neumann entropy is
\bne \label{eq:qmaximin} S(R) = \max_\Sigma \min_{I\subset \Sigma} \left ( \frac{\text{Area}(\del I)}{4G_N} + S_{ren.}(R \cup I) \right )  \ene 
%\cts{Should I define $S_{ren.}$ here? Should also distinguish semiclassical from exact entropy.}

%The maximisation is over time slices containing $R$. 
The quantum maximin prescription \eqref{eq:qmaximin} is equivalent to the island formula, assuming the QFC and a certain stability condition on $I$~\cite{Akers2019}. The prescription minimises over candidate islands on a given time slice, then maximises over time slices. Let us call such a maximising time slice $\Sigma$ a maximin slice. The maximin slice need not be unique, particularly if the island is empty.

\subsection{Proof}

The proof given in~\cite{Bousso2021} is a proof by contradiction. Start by assuming that the two conditions \eqref{eq:cond1a} and \eqref{eq:cond2} are satisfied, and that the true island $I$ is the empty set, i.e. there is no quantum extremal region $I$ that lowers the generalised entropy of $R$.  We will treat the case where $I' \cup R$ is quantum normal first, and consider the quantum anti-normal case separately. 

$I' \cup R$ is quantum normal so the quantum expansion in the inward null directions is initially negative. Consider a time slice $\Sigma$ that contains $R$ and define the representative $\tilde{I}'$ of $I'$ on $\Sigma$ as
\bne \label{eq:tildeI'def} \tilde{I}' := D(I') \cap \Sigma \ene
where $D(I')$ is the domain of dependence of $I'$.  QFC states that the quantum expansion cannot increase along a null congruence, so that the quantum expansion must remain negative along the null congruence from $I'$ to $\tilde{I}'$. The generalised entropy can only decrease under such inward null deformations:
%\cts{Talk about null surface generators, use fancy maths} 
\bne \label{eq:QFC_implication} \begin{split} S_{gen} (\tilde I ' \cup R) &\leq S_{gen} (I' \cup R)\\
& < S_{gen} (R)
\end{split} \ene
The second line follows from \eqref{eq:cond1a}. 

The contradiction appears if we take $\Sigma$ to be a time slice on which the codim-2 quantum maximin surface lives. 
% By definition this is a time slice
%$\Sigma$ on which \cts{Change to notation from quantum maximin paper}
%\bne \max_{\Sigma'} \min_{\mathcal{I} \subseteq \Sigma'} S_{gen} (\mathcal{I} \cup R) = S_{gen} (I \cup R) \ene
%is maximised~\cite{Akers2019}. We assume the existence of such a slice, but it will not in general be unique. \cts{Paragraph needs editing since changed later discussion of quantum maximin.}
From the quantum maximin prescription the true island $I$ gives the lowest generalised entropy on its maximin slice, and since we assumed $I$ to be the empty set we have that
\bne \begin{split} S_{gen}(\tilde I' \cup R) &\geq S_{gen} (I \cup R) \\ &= S_{gen} (R) \end{split} \ene
which directly contradicts \eqref{eq:QFC_implication}. The conclusion is that the island cannot be empty if the two conditions \eqref{eq:cond1a} and \eqref{eq:cond2} are satisfied.

%If we assume the true island is empty, i.e. $I = \varnothing$, then
%\bne S_{gen} (I' \cup R) > S_{gen} (R) \ene
%in contradiction with assumption \eqref{eq:cond1a}.

When $I' \cup R$ is quantum anti-normal the essential steps in the proof are the same, except that the representative $\tilde{I}'$ of $I'$ on a time slice $\Sigma$ is defined as
\bne \tilde{I}' := J(I') \cap \Sigma \ene 
with $J(I')$ all points that can be reached from $I'$ with causal curves~\cite{Bousso2021}. 

\subsection{Loophole} \label{sec:loophole}
%\st{Maximin surface not intersecting domain of dependence of $D$}
The proof in a nutshell is that the conditions imply that every time slice that contains $R$ and intersects $D(I')$ has a non-empty $\tilde{I}'$ for which $S_{gen} (\tilde{I}' \cup R) < S_{gen} (R)$. This rules out that time slice being the maximin slice for a true island that is empty, as the true island must give the lowest generalised entropy on the maximin slice. 

One loophole in the proof is the implicit assumption that the maximin slice of the empty island intersects $D(I')$. If the true island is empty, and its maximin slice does not intersect $D(I')$, then by the definition \eqref{eq:tildeI'def} we have $\tilde I' = \varnothing$, and the specific step that led to a contradiction that fails here is \eqref{eq:QFC_implication}. No contradiction is reached and the proof fails.

\deleted{The QFC does not constrain the generalised entropy of subregions of time slices that do not intersect $D(I')$.}
\added{The QFC does not constrain, in relation to $S_{gen}(I' \cup R)$, the generalised entropy of subregions $X\cup R$, with $X \neq \varnothing$, on time slices that do not intersect $D(I')$.}
QFC is defined on null congruences, and by definition the integral curves in a congruence are non-intersecting, so congruences end at caustics. QFC states that the quantum expansion cannot increase along a null congruence, but the tips of $D(I')$ are caustics. 
\added{The QFC \textit{does} constrain $S_{gen.}(X \cup R)$ for $X = \varnothing$ if we apply it to a quantum normal $I' \cup R$ and integrate it over the whole of either of $\del D^\pm (I')$:}
\bne \label{eq:contra} S_{gen.} (R) \leq S_{gen.} (I' \cup R) \ene
\added{Thus there is a direct tension between island detectors and the QFC, if there exists an empty representative $\tilde{I}' = \varnothing$, because~\eqref{eq:contra} contradicts~\eqref{eq:cond1a}. This is a loophole counterargument that we will return to in Sec.~\ref{sec:Quantum_Bousso_bound}.}

\deleted{This is related to another loophole in the proof, which is the implicit assumption that the null congruence to which the QFC is applied does not end before the past or future tips of $D(I')$ are reached. $\del D^{\pm}(I')$ will not always be a null congruence, because of the integral curves of null generators from $\del I'$ reaching caustics. Equation~\eqref{eq:QFC_implication} in the proof is not implied by QFC if $\tilde I'$ is on a time slice that intersects $D(I')$ but not the null congruence to which QFC is applied. }

\deleted{This second loophole is important, but it is difficult to make general statements about in which set-ups caustics will form and the null congruence end prematurely, so there is little more we can say without studying specific models. By contrast, in general set-ups there are time slices that do not intersect $D(I')$, so in what follows we will focus on the first, `empty representative' loophole.}

\begin{figure}
     \centering
     \includegraphics[width=0.65\textwidth]{Figs/Qantinormal.png}
     \caption{The Penrose diagram of a small AdS black hole that forms from collapse then evaporates away to nothing. $\Sigma$ is a time slice. $I'$ is a region whose representative $\tilde{I'} := \Sigma \cap J(I') = \varnothing$, for which the step~\eqref{eq:QFC_implication} of the proof in~\cite{Bousso2021} is incorrect. This is an example of the causal structure needed for $\tilde{I}'$ to be empty in the case of $I' \cup R$ being quantum anti-normal.}
     \label{fig:antiqnormal}
\end{figure}

%This is an example where the mirage $I' \cup R$ can be anti-quantum normal, while in the other examples the region was quantum normal. 

When $I' \cup R$ is quantum normal any time slice that does not intersect $D(I')$ gives an empty representative $\tilde{I}'$. To have an empty representative when $I' \cup R$ is quantum anti-normal there must exist a time slice that does not intersect $J(I')$. This is not possible in globally hyperbolic spacetimes, because all causal curves from $I'$ will intersect all Cauchy slices. Recall however that quantum maximin is applicable not only to globally hyperbolic spacetimes, but also spacetimes that are stably causal and not globally hyperbolic. This allows for time slices on which $\tilde{I}'$ is empty. An example of a stably causal spacetime with an $I'$ and a time slice $\Sigma$ on which the representative $\tilde{I}'$ is empty is shown in Fig.~\ref{fig:antiqnormal}. Singularities are barriers to causal curves from $I'$. This figure is only an example of the causal structure needed for an island mirage when $I' \cup R$ is quantum anti-normal - we have not shown nor do we claim that the Island Finder conditions are satisfied in this case. 

Whether $I' \cup R$ is quantum normal or anti-normal, there are time slices on which the representative $\tilde I' = \varnothing$, and $S_{gen} (\tilde I' \cup R) \nleqslant S_{gen} (I' \cup R)$. No contradiction is reached, so a true island that is empty is not ruled out. The island finder conditions are insufficient to establish existence of islands, and the conditions may be satisfied without $R$ having a non-empty island. The conditions are sufficient only to establish nonexistence of empty islands whose maximin slice intersects the domain of dependence of $I'$ (when $I' \cup R$ is quantum normal). The concommitant upper entropy bound given in~\cite{Bousso2021} is also not proven if the proof that the conditions are sufficient for a non-empty island has a loophole.

%\st{Notes on QFC and null congruences}
%
%[1510.02099] pg 9: The QFC is defined on null hypersurfaces generated by null congruences
%
%[Carroll] appendix F, by definition the geodesics in a congruence does not intersect each other. "A congruence is a set of curves in an open region of spacetime such that every point in the region lies on precisely one curve".
%
%I need to be careful about applying GSL to my counterexample; I know GSL is violated by evaporating black holes.
%
%[Wikipedia] also says that the integral curves in a congruence are non-intersecting.
%
%QFC says that the quantum expansion cannot increase along a null congruence, but congruences end at caustics. 
%
%The quantum Bousso bound is defined in terms of a quantum light sheet, a null hypersurface on which the quantum expansion is every nonpositive.
%
%If we shoot lightrays from $I'$, and there are caustics, then $L$ is not a null congruence and QFC does not guarantee that $L$ is a quantum light sheet to which the quantum Bousso bound applies.
%
%I'm 90\% sure that QFC only applies to null congruences, i.e. the quantum Bousso bound which is the integrated version only applies up to caustics. I don't have much intuition however whether caustics will generically form. If there's a quantum Raychaudhuri equation, then by the usual argument the quantum expansion would diverge after finite affine parameter. If the quantum expansion is dominated by the classical term, and the null energy condition is not violated by quantum corrections, then the classical Raychaudhuri equation suffices.


%\todo[size=\small]{Why couldn't I interpret this as a proof by contradiction that the empty island slice has to intersect $D(I')$? Why specifically does QFC fail past the tip of the domain of dependence?}

\section{Loophole counterarguments and island puzzles} \label{sec:Puzzles}

\subsection{Time slices for empty maximin surfaces}

The first objection one may have to the loophole is that if the true island is empty, then is it not irrelevant which time slice we use when calculating the entropy? Can we not always choose a time slice with a non-empty intersection with $D(I')$? The loophole is indeed closed if all empty islands have maximin slices which intersect $D(I')$. As a reminder, we defined maximin slices to be the maximising subset of time slices in the quantum maximin prescription~\eqref{eq:qmaximin}.

Now it is true that when the island is empty the von Neumann entropy of $R \cup I$ is the same for all time slices that contain $R$, and that in general a subset of those time slices will intersect $D(I')$. What is not true is that all such time slices
% which contain $R$ on which trivially
%\bne S_{gen} ((I=\varnothing) \cup R) = S_{gen} (R) \ene
are maximin slices for the empty island. If the true island is empty then by maximin there is certainly at least one time slice on which the empty-island entropy is the minimal entropy, but there is no guarantee that the no-island entropy is the minimal entropy on every time slice containing $R$.
% \cts{Give reminder of what a maximin slice is.}

%By definition, maximin slices for $I \cup R$ for general $I$ are the set of time slices $\Sigma$ on which
%\bne \max_{\Sigma} \min_{\mathcal{I}\subset \Sigma} S_{gen}(\mathcal{I} \cup R) = S_{gen} (I \cup R). \ene
Assuming that the island is empty, a time slice fails to be a maximin slice when on that slice there is a subregion $\mathcal{X}$  for which 
\bne \label{eq:noislandcond} \frac{\text{Area}(\del \mathcal{X})}{4G_N} + S_{ren.}(R \cup \mathcal{X})< S_{ren.} (R)\ene
%\bne S_{gen} (\mathcal{I} \cup R) < S_{gen} (R) \ene
As an example of this, consider a time slice that contains a subregion $\mathcal{X}$ whose boundary has zero area, i.e. whose components are null hypersurfaces, and for which $\mathcal{X}$ is entangled with $R$ such that $S_{ren.}(\mathcal{X} \cup R) < S_{ren.} (R)$. Similar to the setups considered in~\cite{Hartman2020}, we could take a thermofield double state on AdS and a non-gravitating bath, take $R$ to be the whole of the bath, and $\mathcal{X}$ to be a ball-shaped region in the AdS whose boundary has been deformed into zigzag union of many null hypersurfaces. The time slices that contain this $\mathcal{X}$ (and $R$) fail to be maximin slices for the empty island.
% because they contains an $\mathcal{I}$ with a lower generalised entropy than the empty island.   
\cts{Think carefully about when/whether $S_{ren.}(AB) < S_{ren.}(A)$.}

We have argued that for an empty island not every time slice (containing $R$) is necessarily a maximin slice, but for the loophole to work we need something stronger: that no time slice that intersects $D(I')$ is a maximin slice when the island is empty. This requires that on all such time slices with $\Sigma \cap D(I') \neq \varnothing$ there exists an $\mathcal X$ with $\mathcal X \cup R \subseteq \Sigma$ 
such that~\eqref{eq:noislandcond} is satisfied.
%\bne  \frac{\text{Area}(\del \mathcal{I})}{4G_N} + S_{ren.} (\mathcal{I} \cup R)  < S_{ren.} (R) \ene
Island Finder~\cite{Bousso2021} showed that~\eqref{eq:noislandcond} is satisfied by the representative $\tilde{I}'$ as defined in~\eqref{eq:tildeI'def}. Thus an $I'$ that satisfies conditions~\eqref{eq:cond1a} and~\eqref{eq:cond2} only rules out empty islands whose maximin slices intersect $D(I')$; it is not sufficient for the existence of a non-empty island.

\subsection{Quantum Bousso bound puzzle} \label{sec:Quantum_Bousso_bound}

Another objection to the loophole is the apparent contradiction with the quantum Bousso bound when there a time slice that does not intersect $D(I')$, with $I' \cup R$ satisfying conditions~\eqref{eq:cond1a} and~\eqref{eq:cond2}. 

The version of the quantum Bousso bound that we use is that which is implied by the QFC~\cite{Bousso2015}. There is a different version of the quantum Bousso bound that has been proven using positivity of relative entropy~\cite{Casini2008, Bousso2014}. In that version of the bound the entanglement entropy is regulated by vacuum subtraction, and vacuum-subtracted entanglement entropy is inequivalent to the renormalised entropy used in $S_{gen.}$~\cite{Bousso2014, Bousso2015}. We need to use the version implied by the QFC, because it is that version that constrains the renormalised and generalised entropies that appear in the island formula.

The quantum Bousso bound applies to quantum light sheets, which are null hypersurfaces on which the quantum expansion is every nonpositive. Since $I'\cup R$ is quantum normal by assumption, QFC implies that the past and future boundaries of the domain of dependence $D(I')$ are quantum light sheets, and the quantum Bousso bound states that~\cite{Strominger2003}
\bne \label{eq:qBoussoBound} S_{gen.}(I' \cup R) \geq S_{gen.} (\tilde I ' \cup R)
\ene
%\bne \label{eq:qBoussoBound} S_{ren.}(I' \cup R) + \frac{\text{Area}(\del I'))}{4G_N} \geq S_{ren.} (\tilde I ' \cup R) + \frac{\text{Area}(\del \tilde{I}'))}{4G_N} \ene
where as before, $\tilde{I}' := \Sigma \cap D(I')$. 
%\added{We do not assume that the time slice intersects $D(I')$; if the representative $\tilde{I}'$ is empty then the quantum Bousso bound states that $S_{gen.} (R) \leq S_{gen.}(I' \cup R) $.}

We assume that the spacetime is stably causal, which is the same causality condition assumed in the quantum maximin prescription, such that there is a global time function $t$. With a reparametrisation of $t$ we can take the $t=0$ slice to contain $I'$, and the $t=\pm 1$ slices to intersect $D(I')$ at the future and past tips of $D(I')$. \cts{Assuming a reparametrisation such that there is a level set that contains $I'$ seems stronger than stable causality.} In the limit of taking time slices to the tips of the domain of dependence, $|t| \to 1^-$, we have $\text{Area}(\del \tilde I') \to 0$. If we assume also that as we deform towards either tip of $D(I')$ that 
\bne \label{eq:dubious_limit} \lim_{|t| \to 1^-} S_{ren.}(\tilde{I}' \cup R ) = S_{ren.} (R) \ene
then the \added{$|t| \to 1^-$ limit} of the quantum Bousso bound~\eqref{eq:qBoussoBound} implies
\bne S_{gen.}(I' \cup R) \geq S_{gen.} (R), \label{eq:qbcont} \ene
%\bne S_{ren.}(I' \cup R) + \frac{\text{Area}(\del I')}{4G_N} \geq S_{ren.} (R), \ene
which directly contradicts \eqref{eq:cond1a}. \added{Thus as we deform $\tilde {I}'$ towards either tip of $D(I')$ then there is a direct tension with the quantum Bousso bound even before the exact point of $\tilde{I}' = \varnothing$ is reached, if~\eqref{eq:dubious_limit} is true.}  

\begin{figure}
     \centering
     \includegraphics[width=0.7\textwidth]{Figs/IslandInBH.png}
     \caption{An evaporating black hole in asymptotically Minkowski spacetime. Radiation region $R$ has an island $I$. Also shown are the boundary of the causal past of $R$, the boundary of the past domain of dependence of $I$, and an earlier time representative $\tilde I$ of $I$. This is a counterexample to the claim that the existence of a region satisfying conditions~\eqref{eq:cond1a} and~\eqref{eq:cond2} implies the existence of singularities which remove both tips of the domain of dependence.}
     \label{fig:BHqboussobound}
\end{figure}

Islands, not just island detector regions, are in direct tension with the quantum Bousso bound by the same chain of logic. Recall that islands are special cases of island detector regions in the sense that they satisfy~\eqref{eq:cond1a} and saturate the bounds~\eqref{eq:cond2}. As an example, consider an evaporating black hole in asymptotically flat spacetime, as depicted in Fig.~\ref{fig:BHqboussobound}. The radiation region $R$ is taken to be sufficiently large to have an island $I$ in the black hole interior. The representative $\tilde{I} := \Sigma \cap D(I)$ on an earlier time slice satisfies
\bne \label{eq:35} S_{gen} (\tilde{I} \cup R) \leq S_{gen} (I \cup R) \ene
by the quantum Bousso bound and QFC. Again we reach the apparent contradiction $S_{gen} (R) < S_{gen} (R)$ when we reach the past tip of $D(I)$ and $\tilde{I} = \varnothing$. 

This contradiction needs to be resolved whenever conditions \eqref{eq:cond1a} and \eqref{eq:cond2} are satisfied, whether or not there is a non-empty island in the spacetime. Here are several ways the contradiction can be evaded:

\vspace{3mm}
1. $D(I')$ is truncated by curvature singularities.
\vspace{3mm}

One possible resolution is that when conditions~\eqref{eq:cond1a} and~\eqref{eq:cond2} are satisfied there are always singularities to the future and past of $I'$ that remove the tips of $D(I')$, so that the problematic limit~\eqref{eq:dubious_limit} does not exist. In some spacetimes this would force all time slices to intersect $D(I')$, so there would be no empty representatives $\tilde{I}'$ and the loophole would be closed
%\footnote{In section~\ref{sec:LongWormhole} we will see a spacetime where $D(I')$ is truncated by singularities yet there are time slices which do not intersect $D(I')$.}
. If truncation by singularities were the only resolution to the tension with the quantum Bousso bound then conditions \eqref{eq:cond1a} and \eqref{eq:cond2} would also be sufficient conditions for the existence of singularities\footnote{This was also pointed out in~\cite{BoussoTalk}.}.

%Naively then the quantum Bousso bound rules out any $I'$ satisfying the island finder conditions \eqref{eq:cond1a} and \eqref{eq:cond2}. The QFC in combination with quantum normalcy imply that the generalised entropy of $\tilde{I}' \cup R$ cannot increase as $\tilde{I}'$ is deformed towards the tip of $D(I')$, it must increase at some point in order to reach $S_{gen} (R)$ at the limit point where $\tilde{I}' = \varnothing$. If we want QFC and the quantum Bousso bounds to be true, and for true islands to exist, then since true islands satisfy the island finder conditions we need to resolution to the above tension. 

%\begin{figure}
%     \centering
%     \includegraphics[width=0.5\textwidth]{Figs/QuantumSingularityCounterexample.png}
%     \caption{An evaporating black hole in asmyptotically flat spacetime. This is a counterexample to the claim that an $I' \cup R$ which is hyperentropic and quantum normal, in the sense of satisfying \eqref{eq:cond1a} and \eqref{eq:cond2}, has a singularity to the past and future of $I'$, which prevents one from taking the limit $\text{Vol}(\tilde{I}') \to 0$.}
%     \label{fig:BHqboussobound}
%\end{figure}



%QFC implies GSL if one assumes the quantum expansion at the asymptotic future of the causal horizon vanishes~\cite{1506.02669}.

% More generally, if a null Cauchy splitting surface has future quantum expansion is positive, it is positive everywhere to the past too.

There is a simple counterexample to the claim that there must be singularities, which is the evaporating black hole of Fig.~\ref{fig:BHqboussobound}. While there is a singularity to the future of $I$ there is none to the past. Curvature singularities may resolve the tension with the quantum Bousso bound in some cases, but there must be a different resolution for the past domain of dependence of an island in an evaporating black hole.

%\comment{Replace with simpler island example.} There is a simple counterexample to the claim that there must be singularities. Consider an evaporating black hole in asymptotically flat spacetime, as depicted in Fig~\ref{fig:BHqboussobound}. The radiation region $R$ is taken to be sufficiently large to contain radiation emitted after the Page time, and we consider an $I'$ that is the interior of the black hole on a late time slice. Condition \eqref{eq:cond1a} is satisfied because the event horizon area is small and there is a large amount of entanglement between the black hole interior and the radiation subregion of $\mathscr{I}^+$. 
%
%The second condition, that $I' \cup R$ is quantum normal, is implied by the generalised second law (GSL). \cts{What about the inward null direction?} The GSL states that the generalised entropy of
%\bne  \Sigma \cap I^- (W_{fut}) = (\tilde{I}' \cup R)^c \ene
%cannot increase as the time slice moves towards the past \cts{rephrase}~\cite{Jacobson2003}. The quantum Bousso bound implies the GSL if one assumes the quantum expansion at the asymptotic future of the causal horizon vanishes~\cite{Bousso2015}. $I^-(W_{fut})$ is the chronological past of a future-directed timelike curve of infinite proper length, which for this evaporating black hole setup we take to end at future timelike infinity $i^+$. $\Sigma \cap I^- (W_{fut})$ does not include $R$ because unlike the causal past $J^-$ the chronological past $I^-$ does not include null-separated points. Assuming the quantum state on $\Sigma$ is pure, then 
%\bne S_{gen.} ((\tilde{I}' \cup R)^c) = S_{gen.} (\tilde{I}\cup R) + S_{c.t.} (\del R) \label{eq:SgenequalsSgen}\ene
%The UV divergent counterterms associated to the boundary of $R$ are irrelevant for quantum normality because they are constant under deformations along null generators of $D(I')$. Using \eqref{eq:SgenequalsSgen} the GSL thus also implies that
%\bne \frac{d}{dt} S_{gen} (R\cup\tilde{I}') \leq 0. \ene
%where $t$ parametrises the time slice that defines $\tilde{I}'$. We thus reach the same contradiction as when applying the quantum Bousso bound in the limit $t \to 1^-$. We start with an $I' \cup R$ that satisfies~\eqref{eq:cond1a} and so has lower generalised entropy that $R$. GSL implies that $S_{gen}(\tilde{I}' \cup R)$ can only decrease as we move the time slice towards the past, but then GSL seems to be violated at some point since $S_{gen}(\tilde{I}'\cup R)$ must increase back up to $S_{gen} (R)$ when $\Sigma$ reaches the past tip of the event horizon and $\tilde{I}'$ becomes the empty set.
%
%%This, in combination with GSL, implies the contradiction $S_{gen} (R) < S_{gen} (R)$ when $\Sigma$ reaches the past tip of the event horizon and $\tilde{I}' = \varnothing$. 
%The aforementioned possible resolution that there is always a singularity that prevents one from reaching the tips of the domain of dependence where $\tilde{I}' = \varnothing$ is not true in this set up: while there is a singularity to the future of $I'$ there is none to the past.

\vspace{3mm}
2. $\del D^\pm (I')$ fails to be a quantum lightsheet\deleted{ due to caustics}.
\vspace{3mm}

Another way in which the quantum Bousso bound can fail to be applicable is if only a subset of $\del D(I')$ is a quantum lightsheet. Classical lightsheets, with everyhere negative classical expansion, can terminate due to the formation of caustics where $\theta \to -\infty$. Assuming the null energy condition is satisfied, the Raychaudhuri equation implies that caustics form, and so lightsheets \added{possibly} end after finite affine parameter whenever the classical expansion is negative~\cite{Raychaudhuri:1953yv}. The truncation of lightsheets underlied an argument against the critically illuminated black hole counterexample to the classical Bousso bound~\cite{Lowe1999}\footnote{Angular caustics also save the classical Bousso bound when there is anisotropy in a collapsing shell, which increases its entropy but also decreases the extent of the lightsheet~\cite{Bousso1999}. The critically illuminated black hole does not violate the vacuum-subtracted entropy version of the quantum Bousso bound~\cite{Bousso2014}. Vacuum-subtracted von Neumann entropy is not classical, in that it is not generally additive, so the resolution in~\cite{Bousso2014} is not applicable to the classical Bousso bound.}. 
\added{This argument against Lowe's counterexample was later realised to be flawed; $\del D(I')$ does \textit{not} fail to be a classical lightsheet, despite null geodesics in the congruence orthogonal to $\del I'$ meeting other null geodesics and caustics~\cite{Akers2018}}\footnote{Thanks to Raphael Bousso for discussions on this point.}.

\added{What is true for classical lightsheets is not necessarily true for quantum lightsheets. If quantum lightsheets can end after finite affine parameter, perhaps because of a quantum analogue of a classical caustic, $\Theta \to -\infty$, then the quantum Bousso bound can fail to apply once the null generator deforming $\del \tilde{I}'$ reaches that point.} 
\deleted{If it is true that quantum lightsheets end after finite affine parameter, because of a quantum analogue of a classical caustic, $\Theta \to -\infty$, then the quantum Bousso bound can fail to apply once the null generator deforming $\del \tilde{I}'$ reaches that point.} 
We will return to the question of whether $\Theta$ must diverge to $-\infty$ in finite affine parameter on a quantum lightsheet in the discussion section~\ref{sec:Discussion}. 

%\cts{Since the QFC only applies to null congruences, then the formation of caustics is another loophole in the proof.}

\deleted{It does not seem likely that caustics can resolve the tension with the quantum Bousso bound for islands for \textit{all} evaporating black hole set ups, though it may for critically illuminated black holes, as the formation of caustics is non-universal in that it depends on the stress-energy tensor of the infalling matter.}
\added{It does not seem likely that $\del D(I')$ failing to be a quantum lightsheet can resolve the tension with the quantum Bousso bound for islands for \textit{all} evaporating black hole set ups, as the truncation of quantum lightsheets seems non-universal in that it depends on the state, and in particular the stress-energy tensor of the infalling matter.}
We would like a more robust resolution for islands in evaporating black holes. 



%I'm 90\% sure that QFC only applies to null congruences, i.e. the quantum Bousso bound which is the integrated version only applies up to caustics. I don't have much intuition however whether caustics will generically form. If there's a quantum Raychaudhuri equation, then by the usual argument the quantum expansion would diverge after finite affine parameter. If the quantum expansion is dominated by the classical term, and the null energy condition is not violated by quantum corrections, then the classical Raychaudhuri equation suffices.

\vspace{3mm}
3. Blueshift effects.
\vspace{3mm}

There is a blueshift effect for Hawking modes near the horizon, and it is tempting to claim that the blueshift of Hawking modes towards the past to super-Planckian energies \replaced{is what leads to violations}{invalidates the use} of the \deleted{semiclassical} quantum Bousso bound\deleted{ and resolves the puzzle}, but we posit that this is not the case.

The renormalised entropy of the island comes from summing over interior Hawking modes below the renormalisation scale. As we propagate the interior modes backwards in time onto the representative region $\tilde I$ the relative blueshift - the ratio of frequencies as measured by free-falling observers that cross the event horizon at different times $v_1$ and $v_2$ - grows exponentially as~\cite{Jacobson2005}
\bne \frac{\omega_1}{\omega_2} \approx e^{(v_2 - v_1)/2r_s} \ene

Towards earlier infalling times the area of the boundary of the representative $\tilde I$ decreases, yet $\tilde I$ continues to be able to contain the same modes which purify the partner modes in $R$ despite the shrinking volume, due to the reduction in wavelength from the blueshift. Once the mode blueshifts to a frequency above the renormalisation scale, whatever that may be, it switches from contributing to $S_{ren.}$ to contributing to the entropy counterterm. Naively the blueshift effect leads to modes with super-Planckian energies that backreact on the geometry, a firewall. 

We are assuming no drama for infalling observers which implicitly fixes the entanglement between interior and exterior outgoing modes such that there is no backreaction~\cite{Almheiri2013}. Altering the state by removing a single Hawking mode from the island \textit{would} lead to a firewall in the past, which would  \replaced{lead to a curvature singularity}{invalidate the use of the quantum Bousso bound}.

To summarise, blueshift effects do not seem to resolve the puzzle because with the no drama assumption there are no curvature singularities in the neighbourhood of the past tip of $D(I)$ that could \replaced{explain the violation}{invalidate the use} of the quantum Bousso bound. 
  
%A related effect is the blue-shift of Hawking modes near the horizon. The renormalised entropy of the island comes from summing over interior Hawking modes below the renormalisation scale. As we propagate the interior modes backwards in time onto the representative region $\tilde I$ the relative blueshift, the ratio of frequencies as measured by free-falling observers that cross the event horizon at different times $v_1$ and $v_2$, is approximately~\cite{Jacobson2005}
%\bne \frac{\omega_1}{\omega_2} \approx e^{(v_2 - v_1)/2r_s} \ene
%Towards earlier infalling times the area of the boundary of the representative $\tilde I$ decreases, yet $\tilde I$ continues to be able to contain the same modes which purify the partner modes in $R$ despite the shrinking volume, due to the reduction in wavelength from the blueshift.  

%Once the mode blueshifts to a frequency above the renormalisation scale, whatever that may be, it switches from contributing to $S_{ren.}$ to contributing to the entropy counterterm $S_{c.t.}$.
%
%An observer standing on the surface of the star as it collapses inside its event horizon should not detect trans-Planckian excitations, assuming no drama~\cite{AMPS}.

%%With a fixed momentum cutoff for a free-falling frame these modes blue-shift past the renormalisation scale and cease to contribute to $S_{ren.}$ of $\tilde I$.
%
%%$\tilde I$ is futher from the curvature singularity than the island $I$.

%Interior and exterior Hawking modes propagate out from a trans-Planckian reservoir close to the event horizon.

%Generalised entropy of states with super-Planckian energy density\footnote{Energy density is of course frame-dependent. If there is a local inertial frame for which the energy density is sub-Planckian then the semiclassical approximation is assumed to be valid~\cite{Wall2010}.} is also not well-defined. 	Consider the evaporating black hole of figure ~\ref{fig:BHqboussobound}. Recall that the island region $I$ satisfies the condition \eqref{eq:cond1b} because at late times the area of the boundary in Planck units is smaller than the reduction in renormalised entropy from including the interior Hawking modes that purify the exterior partner quanta that reach $R$. Now consider deforming $\tilde{I}$ towards the tip of $D^- (I)$, which corresponds time slices progressively further in the past. QFC, or equivalently in this case the quantum Bousso bound, implies the generalised entropy decreases under this deformation, which in turn implies a relation similar to~\eqref{eq:cond1b}:
%\bne \label{eq:puzzle} I(\tilde{I}' , R) > S_{gen.} (\tilde{I}') \ene
%There is a blue-shift effect as we propagate the interior modes on $\tilde I$ backwards in time. The ratio of frequencies of a massless excitations as measured by free-falling observers that cross the event horizon at different times $v_1$ and $v_2$ is approximately~\cite{Jacobson2005}
%\bne \frac{\omega_1}{\omega_2} \approx e^{(v_2 - v_1)/2r_s} \ene
%Interior and exterior Hawking modes propagate out from a trans-Planckian reservoir close to the event horizon.

%%If we neglect quantum fluctuations of the background geometry ...
%
%\comment{To do: Tidy this part up. Get rid of blueshift effects as a resolution to the Bousso bound tension. Move discussion of blueshift effects to Bekenstein bound section. Explain how, even though $\tilde{I}'$ has small proper volume it contains all the Hawking wavepackets that purify $R$, yet it does not backreact strongly.}

\iffalse

We continue deforming until $\tilde{I}'$ is smaller than the local energy and curvature scales, such that the state on $\tilde{I}'$ is approximately that of a reduced vacuum state in flat space. We still take $\tilde{I}'$ to be larger than Planckian scales so we don't have to worry about our semiclassical quantities being well-defined. We reach a puzzle, as how can we have an $O((G_{N}\hbar)^{-1})$ amount of correlation between $\tilde{I}'$ and $R$ when $\tilde{I}'$ is arbitrarily small and approximately in the vacuum state? 

In the evaporating black hole example it is possible because $\tilde{I}'$ is a trans-Planckian reservoir for the interior Hawking wavepackets which purify $R$. As is usual the state 
%on a Cauchy slice long after the black hole is formed from collapse 
is assumed to be in the vacuum as measured by an free infalling observer. The Hawking effect arises from the distinction between free-fall frequency measured by an infalling observer, and the Killing frequency measured by a static obeserver at infinity~\cite{Hawking1975}. As the interior wavepackets of Hawking radiation are propagated backwards in time their Killing frequency is increasingly blueshifted, and sufficiently far in the past the Killing frequency is super-Planckian. Even though $\tilde{I}'$ has small proper volume, it contains all the Hawking wavepackets that purify $R$. 

This resolution claims that the mistake made that led to a tension between the QFC and the existence of islands in the above evaporating black hole example was applying the QFC, and the quantum Bousso bound which the QFC implies, outside of their regimes of validity. For the evaporating black hole, we cannot deform $\tilde{I}'$ all the way to the past tip of $D(I')$ without leaving the regime of validity of QFC, because as we evolve the time slice towards the past the interior modes with support on $\tilde{I}'$ become increasingly blueshifted and eventually reach super-Planckian energies. At these energies the generalised entropy, as a semiclassical quantity, is not even well-defined.

%\comment{Sure that this business with blueshifting does not apply only to the exterior and not interior modes?}

\comment{Understand better how one can have super-Planckian energies as the interior modes are blue-shifted, yet have not backreaction creating a singularity? Rindler decomposition of Minkowski vacuum is a clue.}

\comment{ Want to say something like "The state is still semiclassical in the sense that there is a local inertial frame (the infalling one) in which the energy density is sub-Planckian". Would like to be able to explain why $S(\tilde{I})$ may be separably calculable while $S(R)$ is not.} 

%\comment{Is it the super-Planckian energies or the sub-Planckian length scales of $\tilde{I}'$ that invalidate QFC?}

\fi

\vspace{3mm}
4. The naive result \eqref{eq:dubious_limit} was incorrect due to an unavoidable departure from the semiclassical regime when taking the limit. \added{Violations of the quantum Bousso bound are not necessarily unexpected when the bound is applied outside of the semiclassical regime, such as to regions of sub-Planckian size, or when it is derived from integrating the QFC outside of the semiclassical regime.}
\vspace{3mm}
%\vspace{3mm}
%3.a. Cutoff issues.
%\vspace{3mm}

The last resolution to the contradiction that we'll mention is that the naive result \eqref{eq:dubious_limit} is incorrect, because to take the limit you must consider regions $\tilde I'$ of Planckian size. Generalised entropy is a UV finite and regulator-independent quantity~\cite{Bianchi2012, Cooperman2013}, but it is a semiclassical quantity. The generalised entropy of a region of Planckian size is physically meaningless. It may be  mathematically well-defined if the semiclassical approximation is taken, in the sense that the region can at least be specified if quantum fluctuations of the background geometry are neglected, but the result has no relation to the physical reality. 

A simpler set-up where we can see what goes wrong if we trust formulas for entropy beyond their physical regime of validity is if we consider an interval of length $L$ in a 2d CFT in the vacuum state, whose entanglement entropy is $S(L) = \frac{c}{3} \log (L/\epsilon)$. We can evaluate $S(L)$ for $L< \epsilon$, but it physically meaningless and gives unphysical negative entropy. Just as in~\eqref{eq:dubious_limit}, we can take a naive limit 
\bne \lim_{L \to 0} S(L) = -\infty \ene
and get a result which is not the correct entropy of the empty set, which is zero. There is an important distinction between intervals of infinitesimal and zero length.
%The mistake is that when taking the limit we consider intervals smaller than the cutoff $\epsilon$, which is physically meaningless, and leads to unphysical results such as negative entropy. \cts{There is an important difference between infinitesimal intervals and intervals of zero length.}

\deleted{The conclusion then is that as we shrink $\tilde{I}'$ towards the either tip of $D(\tilde{I}')$ it becomes so small that even if the generalised entropy of $\tilde{I}' \cup R$ can be formally evaluated it is physically meaningless. The semiclassical approximation is invalid and we are beyond the regime of applicability of the QFC. This resolution is how we believe the tension with the quantum Bousso bound is resolved for $\tilde I$ in the evaporating black hole set-up of Fig.~\ref{fig:BHqboussobound}.}

\added{The QFC and the quantum Bousso bound are not physically meaningful when applied to regions of sub-Planckian size, and violations of these bounds are not surprising in these applications as we are outside of the semiclassical regime.}

\added{
This does not immediately resolve the tension with the quantum Bousso bound when it is applied to the whole of $\del D^\pm$ as for an island detector in \eqref{eq:contra}, or for an island as in~\eqref{eq:35}, because in these cases no subregions of sub-Planckian size are directly involved in the bound. 

%This does not immediately resolve the tension with the quantum Bousso bound when it is applied to the whole of $\del D^\pm (I')$, as in \eqref{eq:contra}, because in that case no subregions of sub-Planckian size are directly involved. 

%It does however offer an explanation in cases where quantum Bousso bound is violated. 

We may derive the quantum Bousso bound, as applied in \eqref{eq:contra}, by integrating the QFC over $\del D^\pm (I')$. To do so we must however apply the QFC to sub-Planckian regions, as we integrate towards either tip of $D(I')$, where violations of the QFC are not unexpected. A bound that can be derived by integrating the QFC along a null congruence like $\del D^\pm (I')$ is not always false - that would rule out the GSL~\cite{Bousso2015} - but violations of the bound are not necessarily unexpected since we are applying the QFC outside of the semiclassical regime.}   

% We are not saying that the quantum Bousso bound is always false when it is implied by the QFC only through an application of the QFC outside of the semiclassical regime - that would rule out the generalised second law~\cite{} - only offering an explanation for how the quantum Bousso bound can be violated.

% This does not mean that the quantum Bousso bound is always violated when applied to and arrive at the special case of the quantum Bousso bound given in~\eqref{eq:contra}, but in order to do so we must apply the QFC to sub-Planckian regions.}

%An analogous mistake is made when taking the limit of $\tilde{I}'$ to the tip of $D(I')$: to do so requires considering $\tilde{I}'$ with sub-Planckian length scales, where the generalised entropy is not well-defined.

\subsection{Bekenstein bound puzzle} \label{sec:BekBoundViol}
In the previous subsection we described a tension with the quantum Bousso bound that arises when taking \added{the} infinitesimal limit of the representative of either of an island or a quantum normal island detector\added{, as well as beyond that when the representative is empty}. In this subsection we describe a different puzzle that arises \textit{before} the representative reaches Planckian length scales. 

This puzzle is relevant to the loophole described in section~\ref{sec:loophole}, because just like the Bousso bound puzzle if the only resolution were that curvature singularities truncate $D(I')$ before the tension is reached then it is more difficult for the empty island maximin slice to not intersect $D(I')$, though it is still possible in a stably causal spacetime. 

%\cts{Edit}This is a puzzle that applies to island detectors $I'$, whether or not they are island mirages. Islands $I$ can be thought of as a special case of island detectors, in that they satisfy \eqref{eq:cond1a} and saturate \eqref{eq:cond2}.

Suppose there exists a region $R$ with island detector $I'$, with $I' \cup R$ quantum normal. As we follow null generators along $\del D(I')$, assuming for now that no \deleted{caustics or} curvature singularities are reached, then the QFC implies 
%~\eqref{eq:QFC_implication} which when written in terms of the UV-finite mutual information and generalised entropy is
\bne \begin{split} \label{eq:starteq} S_{gen} (\tilde I' \cup R) &\leq S_{gen} (I' \cup R)\\ &< S_{gen} (R) \end{split} \ene
%\bne \label{eq:starteq} I(\tilde{I}', R) > S_{gen} (\tilde{I}') \ene 
Now in~\cite{Hartman2020} it was argued that any $\tilde{I}'$ which satisfies \replaced{$S_{gen} (\tilde{I}' \cup R) < S_{gen} (R)$}{\eqref{eq:starteq}} also satisfies
\bne \label{eq:BekBekBound} S_{ren.} (\tilde{I}') > \frac{\text{Area}(\del \tilde{I}')}{4G_N} \ene
This means that the representative violates of the Bekenstein area bound. The puzzle arises when $\tilde{I}'$ is parametrically smaller than the local curvature and temperature scale, but still larger than $l_p$. When $\tilde{I}'$ is this small the reduced state on it is approximately the vacuum in flat spacetime, which should not violate the Bekenstein area bound. Since $\tilde I'$ is still larger than the Planck scale we can cannot claim that \replaced{violations of the QFC in~\eqref{eq:starteq} are expected due to the size of the $\tilde{I}'$.}{the QFC was invalid to apply}. \added{Let us assume that the QFC has not been violated and look for other resolutions.}
% on the grounds that $S_{gen}(\tilde I' \cup R)$ is ill-defined.

%Thus we reach another apparent paradox: that $\tilde{I'}$ violates the Bekenstein bound even though it is approximately the vacuum state in flat spacetime. 
\cts{Need reference for this statement about the vacuum state?}    

%\cts{Emphasise that this applies to a super-Planckian size region.}

%Assuming the quantum Bousso bound is valid to apply, i.e. that the quantum lightsheet does not end too early, then this is another puzzle that needs to be resolved even if the spacetime has non-empty islands.

This puzzle needs to be resolved for both islands and island detectors. Let us see how this puzzle manifests itself in a setup with islands. Take an asymptotically flat black hole formed from collapse, see Fig. \ref{fig:BHqboussobound}. By applying QFC to $S_{gen} (I \cup R)$ we get
\bne S_{gen} (\tilde I \cup R) \leq S_{gen} (I \cup R) < S_{gen.} (R) \ene
which again, by the argument given for condition 1 in~\cite{Hartman2020}, implies that the representative $\tilde I$ violates the Bekenstein area bound. If the state on $\tilde I$ is approximately the vacuum then it should not violate the Bekenstein area bound. 

How small $\tilde I$ needs to be before the state on it is approximately the vacuum depends on how the black hole was formed. For stellar collapse $\tilde I$ will be a ball-shaped region in the centre of the star soon after the event horizon formed, and it needs to be smaller than the local temperature scale. For a collapsing shell of radiation, the Vaidya metric, the geometry inside the shell of collapsing radiation is flat and $\tilde I$ can be just inside. In either case, QFC along with the island's quantum extremality property implies that the generalised entropy cannot increase as we deform along null generators from $I$ to $\tilde I$, and we again reach the incorrect conclusion that a ball-shaped region of the vacuum violates the Bekenstein area bound.

\subsubsection{Is Bekenstein's bound necessarily violated?}

To explore this puzzle further we will reproduce the derivation from~\cite{Hartman2020} that
\bne \label{eq:311} S_{gen.} (\mathcal{I} \cup R) < S_{gen.} (R) \implies S_{ren.} (\mathcal{I}) > \frac{\text{Area}(\del I)}{4G_N} \ene
We want to see what approximations were made in the derivation and if those approximations are ever invalid. In~\cite{Hartman2020} the region $\mathcal{I}$ was an island, while for us the region of interest is a representative $\tilde I'$ or $\tilde I$ of an island detector or an island, but since they all satisfy the left hand side of~\eqref{eq:311} the derivation is identical.


%Here are some resolutions:
%\begin{enumerate}
%\item A sufficiently small $\tilde{I}'$ does not exist, say because the past and future tips of $D(I')$ are truncated by singularities. This resolution does not work when $I'$ is a true island region in an flat space evaporating black hole, where there is no singularity to the past.
%\item A small $\tilde{I}'$ exists, but it does not satisfy~\eqref{eq:starteq}. This relation is implied by QFC so it can be untrue if for example the null congruence starting from $\del I'$ ends before reaching $\del \tilde{I'}$, as QFC only applies to null congruences~\cite{Bousso2015}. QFC is also a semiclassical statement, so it does not apply when $\tilde{I}'$ is sub-Planckian in size.
%\item The approximations made to reach \eqref{eq:BekBekBound} are invalid. 
%\end{enumerate}
%The first possibility has known counterexamples, and the second is not satisfactory as it is not clear if or when quantum lightsheets are truncated. 

%\cts{Change $\tilde I'$ to $\mathcal I$.}

\begin{figure}
     \centering
     \includegraphics[width=0.65\textwidth]{Figs/BekBoundProof.png}
     \caption{Representative cartoon of the subregions of time slice $\Sigma$ used in the proof of~\eqref{eq:311}.}
     \label{fig:BekBoundProof}
\end{figure}

The first step is to surround $\mathcal{I}$ with a thin region $C$ of width $\delta \gg l_p$. $C$ shares a boundary with $\mathcal{I}$, such that $\del \mathcal{I} \backslash \del C =\varnothing$, and is spacelike separated from $R$. See Fig.~\ref{fig:BekBoundProof}. Applying monotonicity of mutual information 
\bne I(X,Y\cup Z) \geq I(X,Y), \ene  
to the region complement to $\mathcal{I}\cup C$, which is a superset of $R$, we get 
\bne \begin{split} I(\mathcal{I},(\mathcal{I} \cup C)^c) &\geq I(\mathcal{I} ,R) \\
&> S_{gen} (\mathcal{I})
\end{split} \label{eq:BekBoundWithoutApproximations}\ene
We may rewrite \eqref{eq:BekBoundWithoutApproximations} in the following way, which makes its relation to the Bekenstein area bound transparent,
\bne \label{eq:BekBou} S_{ren.} (\mathcal{I}) - \frac{\text{Area}(\del \mathcal{I})}{4G_N} \geq S_{ren.} (\mathcal{I}) + S_{ren.} (C) - S_{ren.} (\mathcal{I}\cup C). \ene
This is done by expanding the mutual information in terms of the renormalised entropies,
\bne I(\mathcal{I},(\mathcal{I} \cup C)^c) = S_{ren.} (\mathcal{I}) + S_{ren.}(\mathcal{I} \cup C) - S_{ren.}(C), \ene
and rearranging. Note that no approximations have been made so far.

\cts{Is this step valid if $C$ has zero width? I think so as long as the boundaries of $C$ are finite.}

 We reach the conclusion that $\mathcal{I}$ violates the Bekenstein area bound at leading order in $G_N$, 
\bne S_{ren.} (\mathcal{I}) - \frac{\text{Area}(\del \mathcal{I})}{4G_N} \geq 0, \ene
if and only if the right hand side of \eqref{eq:BekBou} is non-negative or $O(G_N^0)$. Conversely, we need
\bne S_{ren.} (\mathcal{I}) + S_{ren.} (C) - S_{ren.} (\mathcal{I} \cup C) \lesssim - O(G_N^{-1}) \ene
in order to avoid the conclusion from \eqref{eq:BekBou} that $\mathcal{I}$ violates the Bekenstein bound. $\mathcal{I}$ and $C$ must have \textit{super}-additive renormalised entropies. Renormalised entropy can have unusual properties for an entropy, including superadditivity and negativity~\cite{Wall2009}, so this is not obviously impossible. The right hand side of~\eqref{eq:BekBou} is \textit{not} the mutual information between $\mathcal{I}$ and $C$, which would be non-negative by subadditivity, because the overlap of the boundaries $\mathcal{I}$ and $C$ is nonzero so there is a leftover piece when the boundary counterterms are added back in.
%, which opens the door of possibility for a region to satisfy \eqref{eq:starteq} without violating the Bekenstein area bound. 

% \eqref{eq:BekBoundWithoutApproximations}, but if we now assume that 
%\bne S_{ren.} (\tilde{I}'\cup C) - S_{ren.} (C) = S_{ren.} (\tilde{I}') + O(G_N^0) \label{eq:approxmade} \ene
%then we reach the Bekenstein area bound violation
%\bne \label{eq:Bek_Bound} S_{ren.} (\tilde{I}') \geq \frac{\text{Area}(\del \tilde{I}')}{4G_N} + O(G_N^0)  \ene

In~\cite{Hartman2020} it is assumed that 
\bne \label{eq:ass1} S_{ren.} (\mathcal{I} \cup C) = S_{ren.}(\mathcal{I}) \ene
and
\bne \label{eq:ass2} S_{ren.} (C) = 0 \ene
up to corrections that are subleading in $G_N^{-1}$ with respect to the area term. This is a reasonable assumption if there exists a time slice on which the state is semiclassical and with a subregion $C$ as defined before, and of width $\delta \gg l_p$ but smaller than the local energy and curvature scales. Then $S_{ren.} (C) \sim \delta^{-(d-2)} \ll G_N^{-1}$. As a first comment, note that one class of examples where there is no $C$ of width $\delta \gg l_p$ is if $\del \mathcal{I}$ and $\del R$ are null-separated. 

\begin{figure}
     \centering
     \includegraphics[width=0.5\textwidth]{Figs/CFigure.png}
     \caption{An evaporating black hole in asymptotically Minkowski spacetime. $C$ surrounds the island $I$, with $\del I \backslash \del C = \varnothing$, and is spacelike separated from $R$. We show that $C$ cannot have both width $\delta \gg l_p$ and trace extrinsic curvature $|K| \ll l_p^{-1}$.}
     \label{fig:CFigure}
\end{figure}

%\cts{New subsubsection?}
\subsubsection{Resolution for islands in evaporating black holes}

As a first step to resolving this subsection's Bekenstein bound related puzzle for an island in an asymptotically flat evaporating black hole, let us re-examine the assumptions~\eqref{eq:ass1} and~\eqref{eq:ass2} for this case. We will show that there is no time slice on which (1) there exists a $C$ of proper width $\delta \gg l_p$ and (2) the curvature scale of the slice, as quantified by its trace extrinsic curvature $K$, is not Planckian. 

The boundaries of $R$ and the island are not null-separated as can be seen in Fig.~\ref{fig:CFigure}, so there \textit{are} time slices with regions $C$ of non-zero width. We are working in the vacuum with respect to an infalling observer, so a natural set of coordinates to use are Painlev\'e-Gullstrand:
\bne ds^2 = - dT^2 + \left (dr + \sqrt{\frac{r_s}{r}}dT \right)^2 + r^2 d\Omega_{d-2}^2 \ene
where $T$ is the proper time for a radially infalling observer initially at rest at infinity. We will do an analysis in a region  perturbatively smaller than the length scale set by the black hole evaporation time, so we neglect the $T$ dependence of $r_s$.

The boundaries of $R$ and $I$ are not null-separated, but on a constant infalling time slice the radial separation of the island and the past lightcone of $R$ is Planckian~\cite{Penington2019}, so on that particular time slice there is no $C$ of spatial width $\delta \gg l_p$.
%, and the approximation $S_{ren}(C) \sim O(1)$ is invalid.
Let $r_{I}$ be the radial position of the boundary of the island on a constant $T$ slice, and $r_{lc}$ the radial position of the intersection of $\del J^- (R)$ with that slice. The reason why the radial separation $(r_{lc} - r_{I})$ must be Planckian is because the gradient of the entropy must be order $O(1/G_N)$ in order to compete with the $O(1)$ gradient of area in the extremality condition
\bne \label{eq:extrem_cond} \del_{r_I} \left ( \frac{\text{Area}(\del I)}{4G_N} + S_{ren.} (I\cup R) \right ) = 0 \ene
The gradient of the area is $\del_{r_I} \text{Area}(\del I) = 8\pi r_I$, while by dimensional analysis the leading divergence of the renormalised entropy is $\del_{r_I} S_{ren.} (I \cup R) \sim 1/(r_{lc} - r_I)$~\cite{Bousso2021}.
%, where $r_{lc}$ is the radial position of the intersection of the past lightcone of $R$ with the $v = v_I$ plane. 
To satisfy the extremality condition \eqref{eq:extrem_cond} thus requires a radial separation $(r_{lc} - r_I) \sim l_p^2 / r_I$. 
%, i.e. the radial separation is Planckian. 

The representative $\tilde I$ at earlier times has an even smaller radial separation from $\del J^- (R)$ than $I$, as the separation of lightrays from the event horizon decreases exponentially in advanced time $v := t + r^*$ towards the past. The maximum width an annular region $C$ surrounding $\tilde I$ on a constant $T$ slice can be is exponentially smaller than for $I$.

Now we discuss other time slicings, as 
%To conclusively resolve the puzzle raised in this subsection, one has to show there is no $C$ on \textit{any} time slice for which the quantities in \eqref{eq:approxmade} are well-defined and the approximation is valid.
we have only shown that there is no $C$ of width $\delta \gg l_p$ on the constant $T$ slice that intersects $\del I$. There are certainly other coordinate systems with time slices on which there is a $C$ of arbitrarily large spatial width, because the null separation of $\del R$ and $\del I$ in the future outward direction is infinite\footnote{It may seem counterintuitive that two nearly-null separated points can have large spatial separation. To get an intuitition for this, consider the spatial geodesic connecting two spacelike separated points in 2d Minkowski. This has proper length $-\Delta x^+ \Delta x^-$ which diverges when $\Delta x^+ \to \infty$ if $\Delta x^-$ is fixed and finite.}. 

$C$ is annular with inner boundary fixed to be $\del I$. For a fixed time slice $\Sigma$ the outer boundary of the maximum width $C$ is $\Sigma \cap \del J^- (R)$, which is as far as $C$ can extend radially while staying spacelike separated from $R$, as shown in Fig.~\ref{fig:CFigure}. 

For fixed $I$ we can increase the maximum width $C$ by deforming the time slice such that $\Sigma \cap \del J^- (R)$ moves towards future null infinity. The cost of increasing the maximum possible width of $C$ by deforming the time slice between $\del I$ and $\del J^- (R)$ towards the null limit is that the slice in this region approaches an annulus-like subregion of a null cone, and null cones have a divergent trace extrinsic curvature.
%Assuming spherical symmetry, $C$ will have two spherical shell boundaries: $\del C = \del \tilde{I}' \cup (\Sigma \cap \del J^- (R))$.  
%For a spherically symmetric timeslice parametrised by $r = r(v)$ we have
%\bne \label{eq:K} K = \frac{1-r'(v)}{r (2r'(v) - f(r))^{1/2}} - \frac{1}{2} \frac{(f(r) - r'(v)) \frac{r_s}{r^2} - 2 r''(v)}{(2r'(v) - f(r))^{3/2}}\ene
%with $f(r) := 1 - r_s / r$. \cts{Double check formula. Check is constant on geodesics.} A time slice which passes through both $\del I$ and $\del R$ allows for a $C$ of arbitrarily large width, but such a slice approaches the null cone. The radius of the null cone follows the null geodesic $2r'(v) = f(r)$, and for which ~\eqref{eq:K} diverges\cts{You haven't convincingly proven that}.

Let us assume that $C$ stays in the neighbourhood ($\Delta T , \Delta r \ll r_s$) of $\del I$, and justify that assumption later. The geometry in this neighbourhood is approximately
\bne \label{eq:approx_geom} ds^2 \approx - dT^2 + d\rho^2 + r_s^2 d\Omega^2_{d-2} \ene
where $\rho$ is the Lema\^itre radial coordinate, which is related to the Schwarzschild radial coordinate by
\bne r(T , \rho) = \left( \frac{3}{2}(\rho - T) \right )^{2/3} r_s^{1/3}. \ene
The radial separations of $\del I$ and $\del J^- (R)$ on a constant $T$ slice are approximately the same in the two coordinate systems
\bne (\rho_{lc} - \rho_i) \approx (r_{lc} - r_i) \sim l_p^2 / r_s \ene
because near $r = r_s$ we have $d\rho \approx dr + dT$.

Now consider a spherically symmetric time slice that intersects $\del I$ and in its neighbourhood has slope $d\rho / dT = \sqrt{1+\alpha^2}$ with $\alpha > 0$ to keep the slice spacelike. The trace extrinsic curvature of this time slice in the neighbourhood of $\del I$ is the same as the surface of a cone in Minkowski spacetime, which is 
\bne \label{eq:Kform} |K| = \frac{1}{\alpha r_s}. \ene
This $K$ diverges in the null limit $\alpha \to 0^+$. Recall that to maximise the possible width of $C$ we want to move the intersection $\Sigma \cap \del J^- (R)$ towards future null infinity, so let us take $1 \gg \alpha \gg l_p / r_s$, which is as close to the null limit as possible while keeping the extrinsic curvature sub-Planckian, $|K| \ll l_p^{-1}$. The slice intersects $\del J^- (R)$ after time 
\bne (\sqrt{1+\alpha^2}-1)\Delta T = (\rho_{lc} - \rho_{I}) \sim \frac{l_p^2}{r_s} \ene
Thus $\Delta T \ll r_s$ given the range of $\alpha$, which justifies the earlier \replaced{approximation~\eqref{eq:approx_geom}}{assumption that the geometry $C$ is embedded in is approximately flat}. The largest possible $C$ on the time slice defined has width given by the radial geodesic length between $\del I$ and the intersection of the slice with $\del J^- (R)$, which is
\bne  \delta = \alpha \Delta T \ll l_p . \ene

%The geometry in a neighbourhood of $\del I$ smaller than the local curvature scale is approximately flat. The radial separation of $\del J^- (R)$ and $\del I$ is of order $l_p^2 /r_s$ on the slice of constant advanced time $v$ that intersects $\del I$, as we argued after~\eqref{eq:extrem_cond}, so we can  approximate the geometry in which $C$ is embedded as Minkowski spacetime with transverse spherical radius $r \approx r_s$, as long as $C$ stays sufficiently close to $\del I$. We consider a spherically symmetric time slice which in the neighbourhood of $\del I$ is a cone of radius $r_s$ and slope $dr/dt_M = \sqrt{1+\alpha^2}$, with $t_M$ the local Minkowski time and with $\alpha > 0$ to keep the slice spacelike. The surface of a spherically symmetric cone in Minkowski spacetime of slope $dr/dt_M = \sqrt{1+\alpha^2}$ has trace extrinsic curvature 
%%\cts{double check}
%\bne \label{eq:Kform} K = \frac{1}{\alpha r_s}, \ene
%at $r = r_s$. This $K$ diverges in the null limit $\alpha \to 0^+$, and to keep $K \ll l_p^{-1}$ we require $\alpha \gg l_p / r_s$. To take $C$ as close to the null limit as we can without getting Planckian extrinsic curvature let us therefore assume that $1 \gg \alpha \gg l_p /r_s$. The slice intersects $\del J^- (R)$ after time 
%\bne (\sqrt{1+\alpha^2}-1)\Delta t_M = (r_{lc} - r_{I}) \sim \frac{l_p^2}{r_s} \ene
%given the initial radial separation. Thus $\Delta t_M \ll r_s$, given the assumed range of $\alpha$, which justifies the earlier assumption that the geometry $C$ is embedded in is approximately flat. The largest possible $C$ on the time slice defined has width
%\bne  \delta = \alpha \Delta t_M \ll l_p . \ene
%%\bne \begin{split} \delta &= \alpha \Delta t_M \\
%%%&\sim \frac{l_p^2}{\alpha r_s} \\
%%& \ll l_p . \end{split} \ene
%%%where $\Delta t_M$ is the Minkowski time elapsed between the intersection of the slice with $\del I$ and $\del J^- (R)$.  

The conclusion is that there is no time slice on which there exists a $C$ with both width $\delta \gg l_p$ and with $|K| \ll l_p^{-1}$. 

We have not shown that islands satisfy the Bekenstein bound, nor do we claim that they do. We have only shown that one particular assumption that is sufficient to show that islands necessarily violate the bound is invalid for the black hole set up of Fig.~\ref{fig:CFigure}. The expectation is that islands violate the Bekenstein bound, while their early time representatives $\tilde{I}$ of size below the local curvature scale, as depicted in Fig.~\ref{fig:BHqboussobound}, do not. In order to avoid the conclusion that the Bekenstein bound is violated we have seen from equation~\eqref{eq:BekBou} that it is necessary for all possible regions $C$ to give strongly ($\sim G_N^{-1}$) superadditive renormalised entropy. 

Let us make some preliminary remarks on how $\tilde I$ and $C$ can have superadditive renormalised entropy, but leave a careful analysis for future work. We have shown that the largest width $C$ can have on a time slice with $|K| \ll l_p^{-1}$ is sub-Planckian for $I$, and it will be exponentially smaller still for $\tilde I$ because of the convergence towards the event horizon of past-directed ingoing light rays. As mentioned in Sec.~\ref{sec:Quantum_Bousso_bound}, semiclassical entropies of regions of sub-Planckian size can be formally evaluated, even if they are not physically meaningful. If $S_{ren.} (\tilde I \cup C) \approx S_{ren.} (\tilde I)$, which seems reasonable as $L_{\tilde I} \gg l_p \gg L_{C}$, then superadditivity requires $S_{ren.}(C) < 0 $. Regions of width smaller than the renormalisation scale can have negative $S_{ren.}$. For example the renormalised vacuum entropy of an interval in a 2d CFT is proportional to $\log (L/\mu)$, with $\mu$ the renormalisation scale, which is negative for $L < \mu$.   For $\tilde I$ the maximum width $C$ can be is exponentially smaller than $l_p$, so $S_{ren.}(C) \lesssim -O(G_N^{-1})$ may be possible. 

%\cts{This argument seems problematic as you've argued earlier that the entropy of regions smaller than the renormalisation scale is not well-defined.}

%so formally at least the renormalised entropy is arbitrarily negative for intervals smaller than the renormalisation scale. If $\tilde I$ has size $\gg l_p$ and $C$ width $\ll l_p$ on the time slice, then $S_{ren.}(C)$ may be arbitrarily negative while $S_{ren.}(\tilde{I}) \approx S_{ren.} (\tilde{I} \cup C)$, and this gives a superadditive entropy. 

% of how $\tilde I$ avoids the Bekenstein bound violation for late

%We can approximate the geometry in which $C$ is embedded as Minkowski spacetime of transverse spherical radius $r \approx r_s$, as both $\del I$ and $\del J^- (R)$ are of Planckian radial separation from the event horizon. 
%\cts{Detail where the approximation breaks down.}
%We consider a spherically symmetric time slice which in the neighbourhood of the event horizon is a cone of radius $r_s$ and slope $dr/dt = (1+\alpha)$, with $\alpha > 0$ to keep the slice spacelike. The maximum width $C$ on the timeslice will be an annulus-like region on the cone with width given the proper distance between $\del I$ and $\del J^- (R)$, which is
%\bne \delta = \alpha \Delta t \sqrt{1+ 2\alpha^{-1}} \ene
%From where the time slice intersects $\del J^- (R)$, we have $\alpha \Delta t \sim l_p^2 / r_s$, so to have $\delta \gg l_p$ we require $\alpha^{-1/2} \gg r_s/l_p$, i.e. $C$ is nearly null as expected, but then from \eqref{eq:Kform} we have that $K \sim l_p^{-1}$ on $C$. \cts{Double check, correct and make precise.}

%\comment{To do: explain how the renormalised entropies of $\tilde I$ and $C$ are super-additive.}
%A time slice close enough to the null limit to allow for a $C$ of width $\delta \gg l_p$ has a Planckian extrinsic curvature. We have established that both the boundary of $\tilde{I}$ and $\Sigma \cap \del J^- (R)$ on a Painlev\'e-Gullstrand time slice are within a Planck length of the event horizon.  The metric close the event horizon, $r-r_s  \ll r_s$, is approximately flat. We have argued that both $\del I$ and the intersection of $\del J^- (R)$ with a constant $T$ slice are within a Planck length of the horizon. We consider a spherically symmetric time slice which in the neighbourhood of the event horizon is a cone of radius $r_s$ and slope $dr/dt = (1+\alpha)$, with $\alpha > 0$ to keep the slice spacelike. Taking the cone towards the null limit $\alpha \to 0^+$ allows for a larger width $C$, but the extrinsic curvature curvature of the cone which is
%\bne K = \frac{1}{r\sqrt{\alpha}} \ene
%diverges in the null limit. To keep the state semiclassical we want $K < l_p^{-1}$ \comment{Why?}. The maximum width $C$ on the timeslice, given by the proper distance between $\del I$ and $\del J^- (R)$, is
%\bne \delta = \alpha \Delta t \sqrt{1+ 2\alpha^{-1}} \ene
%The radial separation of $\del I$ and $\del J^-(R)$ on the $t = 0$ slice is order $l_p^2 / r_s$, to the cone-shaped time slice which   

 

%This is how the apparent implication that the vacuum state violates the Bekenstein bound is resolved for islands in evaporating black holes. 

%On a constant infalling time slice there is no $C$ of spatial width $\delta \gg l_p$ which shares a boundary with $I'$ and is spacelike separated from $R$. 

%It seems problematic to assume that corrections to $S_{ren.}(I \cup C) \approx S_{ren.}(I)$ are $O(1)$ when an $O(1/G_N)$ entropy gradient is needed. \cts{I'm not sure this is true; $S(I \cup R)$ needs a large gradient but not necessarily $S(I)$. Problem might be with $S(C) \approx 0$. }

%Neither islands nor island mirages are ruled out 

% for the island is Planckian, and when deforming the island towards the past tip of $D(I)$, the radial separation of the representative $\tilde I$ and the past lightcone of $R$ decreases exponentially \cts{Can improve this formula}
%\bne \Delta r \sim \exp ((v-u_0)/2r_h) \ene 
%at earlier infalling time $v$. \cts{Define $r_h$. It is the horizon radius? If so, at what time? Or is it simply some function of the initial black hole mass?}

%\cts{Discuss other $C$ which does have $\delta \gg l_p$. There is a $C$ of large spatial width, but there is a large relative boost.}

%\cts{Write something in discussion about the AdS$_2$ examples. Comment on singularities where dilaton vanishes.}
%\comment{Work in progress from this point. Currently I am looking at how this Bekenstein bound paradox is explicitly resolved in JT gravity island examples such as in figure~\ref{fig:IOTH} from \cite{Almheiri2019}:}
%\begin{figure}[h!]
%     \centering
%     \includegraphics[width=0.5\textwidth]{Figs/IOTH.png}
%     \caption{Two sided finite temperature AdS$_2$ black hole coupled to flat space baths. \comment{Figure taken from \cite{Almheiri2019}. How does the representative of the island near the tip of $D(I)$ evade the Bekenstein bound paradox?}}
%     \label{fig:IOTH}
%\end{figure}

%  $\tilde{I}'$ that  satisfy \eqref{eq:QFC_implication} also violate the Bekenstein bound 

\iffalse

\st{Casini qCEB proofs}

A version of the quantum Bousso bound has been proven

There are two ways which Casini-type proofs of the quantum Bousso bound are not applicable: their definitions of entropy are different, they use vacuum-subtracted EE rather than the renormalisation that occurs in generalised entropy. Further, they require the gravitational backreaction to be small.

[https://arxiv.org/pdf/1404.5635.pdf] section 2 shows that the vacuum subtracted entropy on a null interval, which has zero proper length, is "finite and well defined". This is a big problem for your claim that the quantum Bousso bound can be violated by cutoff subtleties. There may be subtlety regarding the difference between vacuum-subtracted EE, or blue-shift effects. 

Vacuum subtracted EE as used in the Casini-type proofs of the quantum Bousso bound are not the same as the renormalised entropies used in generalised entropy. In the latter we just subtract off the leading UV divergences associated with the regions' endpoints. It is easy to check that $I(A:B) = S(A) + S(B) - S(AB)$ can be rewritten with the EE's replaced by their renormalised equivalents, but not their vacuum-subtracted cousins.

[1404.5635] eq (5.2) and discussion around it is about renormalised entropy as used in the generalised entropy, vs vacuum subtracted EE. The takeaway is that they are not the same.

[9911002] pg 2-3: The Bekenstein bound was only ever supposed to be valid for regions where gravity is not the dominant force. This is why black hole interiors and "super-horizon" size regions in FRW cosmologies can violate the Bekenstein bound.

This is consistent with the Casini-type proofs of the Bekenstein and quantum Bousso bounds, which are only valid for states with negligible gravitational backreaction.
\fi

%\comment{To do: Abridge and edit the following section.}

\section{Discussion and future work}\label{sec:Discussion}

In section~\ref{sec:loophole} we pointed out a loophole in the proof given in~\cite{Bousso2021} that their set of conditions are sufficient for the existence of a non-empty island. There are trivial fixes to their set of conditions to make them sufficient, such as an additional condition that all time slices in the spacetime intersect $D(I')$ (or $J(I')$ if $I' \cup R$ is anti-quantum normal rather than quantum normal). This particular additional condition would be undesirably restrictive though: it is satisfied for the recollapsing FRW universe example of~\cite{Hartman2020} but not for evaporating asymptotically AdS or flat Schwarzschild black holes. The bottom line is that we do not have a set of conditions that are (1) sufficient to establish existence of a non-empty island, (2) easier to evaluate than the island formula, and (3) are not overly restrictive in that they fail to detect islands for whole classes of examples where islands are expected or have been explicitly found.

Having pointed out loopholes in the proof, it is natural to look for explicit counterexamples, i.e. island mirages. In appendix~\ref{sec:Mirages} we outline two promising setups which we intend to explore further. There may be examples of island mirages which are very different in character from those suggested as possibilities in the appendix.

In section~\ref{sec:Quantum_Bousso_bound} we mentioned that it is unclear when or whether a quantum lightsheet will truncate after finite affine parameter. Such a result would be important as it limits the applicability of the quantum Bousso bound. Classically, \replaced{null geodesics in a congruence}{lightsheets} with negative expansion \replaced{reach}{end at} caustics after finite affine parameter because of the bound
\bne \label{eq:classfocussing} \theta' \leq -\frac{1}{d-2}\theta^2 \ene
which follows from the Raychaudhuri equation for null geodesic congruences and assuming the null energy condition. Similarly, a negative quantum expansion leads to the `quantum caustic' $\Theta \to -\infty$ after finite affine parameter if the following quantum analogue of~\eqref{eq:classfocussing} is true:
\bne \label{eq:strong-QFC} \frac{\delta}{\delta V(y_2)} \Theta[V;y_1]\leq -\frac{1}{d-2} \Theta (y_1) \Theta (y_2) \delta (y_1 - y_2) \ene
using the notation from~\cite{Bousso2015}.  This conjecture~\eqref{eq:strong-QFC} is stronger than QFC. The on-diagonal components give \eqref{eq:classfocussing} in the classical limit, and give a stronger version of QNEC in the weak gravity limit\footnote{Which in 4d may not be satisfied by a ball-shaped region of the CFT vacuum in Minkowski spacetime. We thank Aron Wall for pointing this out.}. We feel that these conjectures and their implications for the quantum Bousso bound are worth further study.

In section~\ref{sec:BekBoundViol} we described and gave possible resolutions to the Bekenstein bound puzzle, and discussed the puzzle in the context of islands in evaporating black holes in detail. We did not however give a complete argument for how the early time representative $\tilde I$ of the island can avoid the implication that it violates the Bekenstein bound with superadditive renormalised entropies, as it seems it must. For future work we believe that the islands in AdS$_2$ black holes found in~\cite{Almheiri2019} would be a fruitful arena in which to explore this puzzle in explicit calculational detail.





\acknowledgments

I would like to thank Tarek Anous, Raphael Bousso, Sean Colin-Ellerin, Ricardo Espindola, Jackson Fliss, Bahman Najian, Dominik Neuenfeld, Dora Nikolakopoulou, Edgar Shaghoulian, Arvin Shahbazi-Moghaddam, Shan-Ming Ruan, Manus Visser, and Aron Wall for useful discussions. Special thanks to Ben Freivogel for many productive discussions and comments on an earlier draft. This research was supported by the Stichting Nederlandse Wetenschappelijk Onderzoek Instituten (NWO-I). 
%
\appendix
%
\section{Possible island mirages} \label{sec:Mirages}

%\comment{This section contains my proposed counterexamples. None of them are 100\% convincing yet, and I indeed do not expect all of them to work and to make it to the final draft. One convincing example would be sufficient.}

An island mirage is an island detector $I'$ that satisfies conditions~\eqref{eq:cond1a} and~\eqref{eq:cond2}, but for which $R$ has no island. Mirages, if they exist, are counterexamples to the claim that the conditions are sufficient to guarantee that $R$ has an non-empty island somewhere. In this appendix we outline two setups which we were not able to explicitly show have island mirages, but that we feel are candidates worth further study in future work. 
%\cts{This paragraph was wrong, because there are subregions of the time reflection symmetric slice of dS that are hyperentropic, and are q.normal/antinormal by symmetry.} Mirages are false positives, in that they satisfy the Island Finder conditions the radiation region having anywhere an island. There are also false negatives where the conditions are not satisfied but there is an island, such as for a thermal state on $dS_d$ with $d>2$. Examples include islands which are the whole Cauchy slice of a closed universe. As shown in where no subset of a Cauchy slice of a closed universe satisfies the necessary conditions of islands, such as for thermal states on $dS_{d>2}$, but the whole Cauchy slice is an island, as studied in~\cite{Hartman2020, Bousso2022, Espindola2022}. \comment{Sure they're not quantum normal/antinormal?}

%We have explained that to be an island mirage it is necessary that $I'$ be quantum normal and, except for exceptional cases as described in \ref{sec:BekBoundViol}, Bekenstein bound violating. These coincide with the first two of the set of necessary conditions that true islands must satisfy, as given in~\cite{Hartman2020}. In that paper, the authors suggest that these may in fact be sufficient conditions for the existence of an island, based on a lack of examples where the first two conditions are satisfied but the third is not. If this were true then there could be no counterexamples to Island Finder. The third necessary condition from~\cite{Hartman2020}, that regions $G$ which surround and share a boundary with $I'$ must be quantum normal, is not implied by the Island Finder conditions. In fact the opposite is true: quantum normality of $I' \cup R$ implies that $G = (I'\cup R)^c$ is \textit{anti}-quantum normal. An $I'$ that satisfies the conditions~\eqref{eq:cond1a} and~\eqref{eq:cond2} cannot therefore be an island since it does not satisfy one of the necessary conditions of islands. This alone does not mean the conditions in \cite{Bousso2021} are insufficient to prove island existence, as the region $I'$ is a detector of islands rather than an island itself. 



\subsection{Small, critically illuminated, isolated AdS black holes}
%\comment{I've spent the most time on this counterexample, but it may be a dead end. Regions $I'_a$ and $I'_b$ in figure~\ref{fig:CriticallyIlluminated} seem to be ruled out. $I'_c$ might work, but without a clear calculation it is not convincing.}

\begin{figure}
     \centering
     \includegraphics[width=0.8\textwidth]{Figs/CriticallyIlluminated.png}
     \caption{A small, critically illuminated, isolated AdS black hole, and an auxiliary non-gravitating flat spacetime. The radiation shown includes the incoming radiation that keeps the black mass constant, the outgoing Hawking modes, and their interior partners. The ingoing radiation is maximally mixed and purified by the degrees of freedom in $R$ in an auxiliary spacetime.}
     \label{fig:CriticallyIlluminated}
\end{figure}

Here we consider a small isolated AdS black hole that is critically illuminated by mixed state radiation that is purified by an auxiliary non-gravitational system, see Fig.~\ref{fig:CriticallyIlluminated}. Unlike typical island setups, isolated AdS black holes are not coupled to a non-gravitational bath; the asymptotic boundary conditions are reflecting for all time. Critically illuminated black holes have an incoming flux of energy that exactly balances the outgoing flux due to Hawking radiation, such that black hole has constant mass for some finite period of time~\cite{Lowe1999}. In our setup what an external observer sees is the formation of a small AdS black hole from collapse, that is then critically illuminated for some Schwarzschild time $t_{illum.}$, then allowed to evaporate away before the earliest Hawking radiation has had time to reflect and return from the boundary. For AdS black holes to be small enough to evaporate requires a mass less than
\bne \label{eq:massbound} M l_{AdS} \ll \left ( \frac{l_{AdS}}{l_p} \right )^\frac{d^2 -1}{2d -1} \ene
and there are no small AdS$_d$ black holes in $d \leq 3$~\cite{Harlow2016}. 

The parametric separation of time scales is taken to be
\bne t_{p} \ll t_{evap.} \ll t_{illum.} \ll l_{AdS} \ene 
There are several reasons to look for an island mirage in such a setup:

\begin{enumerate}
\item Isolated AdS black holes do not have islands. Every point in the black hole interior lies on a past causal horizon, so GSL rules out any quantum extremal surface~\cite{Jacobson2003, Engelhardt2021}.
\item Critically illuminated black holes are a counterexample to the classical Bousso bound as an arbitrarily large amount of classical entropy can pass through the event horizon while the horizon area stays constant~\cite{Lowe1999}. An arbitrarily large amount of entropy passing through the horizon is helpful for satisfying condition~\eqref{eq:cond1a}.
%Either a violation of the classical Bousso bound, or an evasion of it applying to the whole of $\del D^- (I')$, is necessary for an island mirage. \cts{Is that true?}
\item Small isolated AdS black holes only differ from asymptotically flat black holes, which do have islands, far from the black hole.
%\item As the black hole is allowed to evaporate to nothing, it is reasonable that the same mechanism that creates islands in AdS m
\end{enumerate}

%\cts{Paragraph is jarring.} 
Let us give a few comments about critically illuminated black holes. They are a counterexample to the classical Bousso bound, though not to GSL as the Hawking radiation continually adds to the entropy outside the black hole but it does not cross the lightsheet~\cite{Lowe1999}. The classical Bousso bound is in effect violated by quantum effects in the form of Hawking radiation. The critically illuminated black hole requires fine tuning of the incoming radiation, and microscopic fluctuations in the incoming radiation generally leads to the lightsheet truncating in finite affine parameter due to caustics
%. Since the classical Bousso bound only bounds the entropy flux through the lightsheet, if ingoing null congruences passing though incoming radiation end in caustics then the lightsheet can end before a sufficient amount of entropy has crossed it to violate the bound
~\cite{Flanagan2000, Bousso2000, Bousso2002}. 

%\begin{table}[]
%\centering
%\begin{tabular}{|l|l|l|}
%\hline
%       & $S_{gen} (I' \cup R) < S_{gen}(R)$ & $\pm \Theta_{\pm} (I' \cup R) \geq 0$ \\ \hline
%$I'_a$ & \checkmark                         & $\times$                              \\ \hline
%$I'_b$ & $\times$?                           & \checkmark                            \\ \hline
%$I'_c$ & ?                                  & ?                                     \\ \hline
%\end{tabular}
%\caption{\comment{Summary of regions $I'$ considered in figure~\ref{fig:CriticallyIlluminated} and which of the Island Finder conditions they satisfy. Table won't be in final paper.}}
%\end{table}
Since there can be no islands in isolated AdS black holes, to find an island mirage we just need to find an $I'$ that satisfies the conditions \eqref{eq:cond1a} and \eqref{eq:cond2}. In Fig.~\ref{fig:CriticallyIlluminated} we depict three such candidate $I'$ wich we will now discuss in turn. $I'_a$ satisfies \eqref{eq:cond1a} because it contains all the ingoing radiation so $S_{ren.}(I'_a \cup R) \approx 0$, and the boundary of $\del I'_a$ has small area, but it fails to be quantum normal because it is deep in the classically trapped region. $I'_b \cup R$ is quantum normal by GSL, because $\del I_b'$ lies on both a future and past causal horizon, but it seems likely to fail to satisfy \eqref{eq:cond1a} because it contains all the interior modes. 
%To satisfy \eqref{eq:cond1a} the Page curve of the exterior $G = (I'_b \cup R)^c$ would have to dip below: $S_{S(R)
$I'_c$ is in a sense in between $I'_a$ and $I'_b$, the Goldilocks candidate. The future inward null quantum expansion is non-positive by GSL, but since $\del I'_c$ is inside the event horizon it does not lies on a future causal horizon, so GSL does not apply to the future outward direction. To be quantum normal it must be outside the quantum trapped region, yet it must also contain enough of the incoming radiation to satisfy~\eqref{eq:cond1a}, and it is not clear whether both of these can be simultaneously satisfied.
%   In future work we would like to see if $I'_c$ is an island mirage. 

\subsection{Large AdS black hole with infalling shell}
\begin{figure}
     \centering
     \includegraphics[width=0.99\textwidth]{Figs/NullShellv2.png}
     \caption{Left: A large, isolated AdS black hole formed from a collapsing star. The star is in a mixed state of entropy $S_{star}$ that is purified by $R$ in an auxiliary non-gravitating spacetime. Spherical symmetry is assumed, and $I'$ are ball regions whose boundaries are points in the figure. To be an island mirage $I' \cup R$ must be both outside the quantum apparent horizon, depicted in blue, and the inside the red hyperentropic region. Right: A null shell of pure radiation is added. The null shell has a large entropy gradient in the $x^+$ direction, so as to push the quantum apparent horizon into the red hyperentropic region. }
     \label{fig:NullShell}
\end{figure}

In this example we consider a large, isolated AdS black hole formed from a collapsing star, as shown in Fig. \ref{fig:NullShell}. The star is in a mixed state, with entropy $S_{star}$. The star is purified by $R$ in an auxiliary non-gravitating spacetime.
% such that $S(\text{star}\cup R) \approx 0$. 
As an isolated AdS black hole, GSL again rules out non-trivial quantum extremal surfaces, so there are no islands. We assume spherical symmetry and look for ball-shaped island mirages $I'$, centred on the origin, which satisfy the conditions \eqref{eq:cond1a} and \eqref{eq:cond2}. 

To satisfy \eqref{eq:cond1a} and lower the generalised entropy, the boundary of $I'$ must be in the red region of Fig. \ref{fig:NullShell}. There the surface area of $I'$ is small and $I'$ contains a sufficiently large fraction of the collapsing star, which by construction is entangled with $R$, to lower the generalised entropy. To satisfy \eqref{eq:cond2}, $\del I'$ must be outside the quantum trapped region, whose boundary is the quantum apparent horizon depicted in blue in Fig.~\ref{fig:NullShell}. 

Take a point at the radial origin of the star after it has collapsed inside its Schwarzschild radius. The subset of the future directed null cone that starts from this point to a cut of the lightcone outside the classical trapped region is a lightsheet. If we assume the star's entropy to be extensive, then the classical Bousso bound places an upper limit on the entropy flux through the lightsheet
\bne S(L) \leq \frac{\text{Area}(\del L)}{4G_N}. \ene
\cts{$<$ or $\leq$?} This rules out an $I'$ whose boundary is outside the classically trapped region, and whose classical entropy violates the Bekenstein bound. If we assume that this island detector region violates the Bekenstein bound, i.e. there is a $C$ for which $I'$ and $C$ do not have superadditive renormalised entropies (see section~\ref{sec:BekBoundViol}), then we reach the conclusion that an island detector cannot exist outside the classically trapped region if the entropy of the matter forming the black hole is approximately classical, i.e. extensive. 
\cts{Think carefully about that statement.}

We need to add matter whose entropy is not classical in order to evade the classical Bousso bound, in order to have an $I'$ that is both quantum normal and violates the Bekenstein bound~\cite{Hartman2020}. We would like to `push' the quantum apparent horizon inward, as depicted in the right hand side of Fig. \ref{fig:NullShell} by creating a large positive entropy gradient in the $x^+$ direction, so making $\Theta_+$ positive where it was negative before. 

%\comment{The rest of this subsection is a sketch of ideas about this null shell. The basic idea is shown in the right half of figure~\ref{fig:NullShell}: to try to push the quantum apparent horizon inside the red region. As I argue, such a possibility is not ruled out by QFC, but an explicit calculation would be best.}

We can try adding a thin shell of radiation of thickness $\Delta x^+$. Take the shell to be composed of two concentric layers that are maximally entangled with each other, but with the whole shell in a pure state. A ball shaped region that is concentric with the shell will then have a large entropy gradient as a function of radius at the shell radius. 
%I imagine the shell to add a triangle-wave profile to the vN entropy of $I' \cup R$ which is a function of the $x^+$ position of $\del I'$:
%\bne \Delta S_{I' \cup R} (x^+) \sim S_{peak} \left(1- \frac{2}{\Delta x^+} |x^+ - x_0^+ | \right), \quad \text{ for } \quad |x^+ - x_0^+ | \leq \Delta x^+ / 2 \ene 
%We take $S_{peak} \ll S_{star}$ to try to keep the shell's effect on the red region small.
A positive entropy gradient in the $x^+$ direction has the potential to flip the sign of $\Theta_+$ and so push the quantum apparent horizon in. 

We do have to be careful about the backreaction of the null shell  though and the effect on the classical expansion. One promising sign is that the quantum apparent horizon of the right hand side of Fig. \ref{fig:NullShell} seems not to violate the QFC. The QFC would forbid $\Theta_+ (I' \cup R)$ from changing signs twice as we deform $I'$ in the $x^+$ direction, but here it only changes sign once. A more careful analysis of the interplay between the entropy gradient and backreaction of the null shell is needed before this is a convincing example of an island mirage.

\cts{Add that radiation is free (pressure would stop it collapsing)}

%\cts{Move possible mirages to appendices?}
 
\iffalse
\subsection{Minkowski vacuum}
\begin{figure}
     \centering
     \includegraphics[width=0.3\textwidth]{Figs/Mink.png}
     \caption{Vacuum state in Minkowski spacetime. Spherical symmetry is assumed and $\del I'$ and $\del R$ are null-separated.}
     \label{fig:Mink}
\end{figure}
\comment{This subsection is currently the least developed of the possible island mirages. I give bullet point lists of why there might be an island mirage here, and where I am in showing the conditions \eqref{eq:cond1a} and \eqref{eq:cond2} are satisfied.}

\noindent Why this setup, Fig.~\ref{fig:Mink}, is promising:
\begin{itemize}
\item No islands in Minkowski vacuum. Can argue this from GSL, or in a more handwacy way that there is no info paradox for islands to resolve here. 
\item In many ways is similar to an evaporating flat space black hole (see Fig. \ref{fig:BHqboussobound}), where there is an island, in the following ways: 
\begin{itemize}
\item $R$ collects radiation at null infinity, just Rindler radiation instead of Hawking radiation.
\item $I'$, a region behind a causal horizon, purifies $R$.
\end{itemize}
\item $I'$, as a ball region in Minkowski vacuum, is not Bekenstein bound violating. Since $\del I'$ and $\del R$ are null-separated, there is no region $C$ of small but finite width between them, so a key step that leads from \eqref{eq:cond1a} to $I'$ being Bekenstein bound violating~\cite{Hartman2020} is evaded. 
\end{itemize}
Are all Island Finder conditions satisfied?
\begin{itemize}
\item $\Theta_+ (I' \cup R) \geq 0$ \checkmark \, by GSL
\item $\Theta_- (I' \cup R) \leq 0$ ?
\item $S_{gen}(I' \cup R) < S_{gen} (R)$ ? Stuck on calculating the entanglement entropy of $I' \cup R$. The complement region is a subregion of the null cone.
\end{itemize} 
\fi

\iffalse
\subsection{Long wormhole}\label{sec:LongWormhole}

\comment{My two concerns with this counterexample are that I may be wrong about how the system evolves in time, or that I am missing some subtlety in the argument for the true island being empty. It doesn't seem right that I can have as much entropy in the wormhole as I like without changing the event horizon area. I would like to understand how fast the wormhole collapses, to look for similar examples in the literature, and whether $I'$ will be q. normal or antinormal. I might be worried that if just the ends of the wormholes that close off, and you're left with a disconnected closed Universe where the island lives.}

\begin{figure}
     \centering
     \includegraphics[width=0.99\textwidth]{Figs/LongWormhole.png}
     \caption{Two small AdS black holes connected by a long wormhole. $t_L = t_R = 0$ is the time reflection-symmetric slice in the spacetime with two asymptotically AdS boundaries. The state on this slice is thermal and purified by the auxiliary system $R$ in the non-gravitational spacetime.}
     \label{fig:LongWormhole}
\end{figure}

Our next example is illustrated in Fig.~\ref{fig:LongWormhole}. We construct our setup by specifying a time reflection symmetric state on a Euclidean geometry, then evolve forwards and backwards to complete the spacetime. We specify the spatial geometry to be a long, thin wormhole that connects two asymptotically hyperbolic regions. The geometry has spherical symmetry and a $\mathbbm{Z}_2$ spatial reflection symmetry about the centre of the wormhole. The state is a static thermal gas which is purified by an auxiliary $R$. If we evolve forwards or backwards in Lorentzian time then the wormhole quickly collapses and the AdS regions disconnect. The spacetime is not globally hyperbolic but it is stably causal which is sufficient for the quantum maximin prescription~\cite{Akers2019}. 

 \cts{Look at literature on bags-of-gold for small AdS black holes.}

To an external observer in either of the AdS regions what is seen is a shell of radiation that collapses to form a small black hole (of mass~\eqref{eq:massbound}) which evaporates to nothing, all in a time-reflection symmetric way. \cts{Comment on initial value problem and naked singularities?}.

To establish that there is an island mirage we need to find an $I'$ that satisfies the two conditions~\eqref{eq:cond1a} and~\eqref{eq:cond2}, and argue that there is no non-empty island. Take $I'$ to be a $\mathbbm{Z}_2$-symmetric subregion of the wormhole on the time reflection symmetric slice, and whose boundaries $\del I'$ are causally disconnected from the AdS boundaries, as depicted in Fig.~\ref{fig:LongWormhole}. 

Condition~\eqref{eq:cond2} is trivially satisfied, as by time reflection symmetry $I' \cup R$ must be either quantum normal and anti-normal. It cannot be trapped or anti-trapped.

The state on the time reflection-symmetric slice is a mixed thermal state with an extensive entropy, and the state is purified by $R$. This gives us
\bne S_{gen.} (I' \cup R) \approx S_{gen.} (R) + \frac{\text{Area}(\del I')}{4 G_N} - s_{th} \text{Vol}(I') \ene
where $s_{th}$ is the thermal entropy density. If we fix the width of the wormhole, and the fraction of the length that $\tilde{I}'$ occupies, then for a sufficiently long wormhole the condition \eqref{eq:cond1a} is satisfied. 

\cts{Argument is not rock solid.} To establish that $I'$ is an island mirage for $R$, we last thing we have to show is that the true island for $R$ \textit{is} empty. First note that the exact von Neumann entropy of $R$ as computed by the island formula should not depend on where the time slice meets the AdS boundaries. This is because $R$ is decoupled from the asymptotically AdS gravitational system. 
Now consider a time slice that intersects the AdS boundaries long after the two small black holes have evaporated. The spatial geometry is two disconnected pure AdS regions, and the state on each region is mixed Hawking radiation from the evaporation of their respective black holes. Pure AdS with thermal radiation has no non-empty islands~\cite{}, so the island formula tells that the exact entropy of $R$ is the total thermal entropy of the Hawking radiation on the two disconnected regions of the slice. The semiclassical matter entropy of every time slice of the AdS system is the same. There cannot therefore be a non-empty island for any choice of AdS boundary time. If there were a non-empty island at an earlier time then the exact entropy of $R$ would depend on AdS boundary time, which it cannot.  

This setup has some similarities to the bag-of-gold example considered in \cite{Bousso2021}, where they consider a highly entropic region behind a black hole event horizon in an asymptotically AdS spacetime. In that case there is a non-empty island. The key difference from that example is that the black holes in this setup are small and can evaporate away so there exists time slices which do not intersect $D(I')$ or $J(I')$, which crucially allows for empty representatives $\tilde{I}'$. 
\fi

\bibliography{references}
\bibliographystyle{JHEP}





\end{document}
\providecommand{\href}[2]{#2}\begingroup\raggedright\begin{thebibliography}{10}

\bibitem{Bousso2021}
R.~Bousso and A.~Shahbazi-Moghaddam, {\it Island finder and entropy bound},
  {\em Phys. Rev. D 103, 106005 (2021)} (Jan., 2021)
  [\href{http://arxiv.org/abs/2101.11648}{{\tt arXiv:2101.11648}}].

\bibitem{Hawking1975}
S.~W. Hawking, {\it Particle creation by black holes},  {\em Commun. Math.
  Phys.} {\bf 43} (1975) 199--220. [Erratum: Commun.Math.Phys. 46, 206 (1976)].

\bibitem{Hawking1976}
S.~W. Hawking, {\it Breakdown of predictability in gravitational collapse},
  {\em Phys. Rev. D} {\bf 14} (1976) 2460--2473.

\bibitem{Page1993}
D.~N. Page, {\it Information in black hole radiation},  {\em Phys. Rev. Lett.}
  {\bf 71} (1993) 3743--3746, [\href{http://arxiv.org/abs/hep-th/9306083}{{\tt
  hep-th/9306083}}].

\bibitem{Ryu_2006}
S.~Ryu and T.~Takayanagi, {\it Holographic derivation of entanglement entropy
  from the anti–de sitter space/conformal field theory correspondence},  {\em
  Physical Review Letters} {\bf 96} (May, 2006).

\bibitem{Freedman2017}
M.~Freedman and M.~Headrick, {\it Bit threads and holographic entanglement},
  {\em Commun. Math. Phys.} {\bf 352} (2017), no.~1 407--438,
  [\href{http://arxiv.org/abs/1604.00354}{{\tt arXiv:1604.00354}}].

\bibitem{Hubeny2007}
V.~E. Hubeny, M.~Rangamani, and T.~Takayanagi, {\it A covariant holographic
  entanglement entropy proposal},  {\em JHEP} {\bf 07} (2007) 062,
  [\href{http://arxiv.org/abs/0705.0016}{{\tt arXiv:0705.0016}}].

\bibitem{Engelhardt2014}
N.~Engelhardt and A.~C. Wall, {\it Quantum extremal surfaces: Holographic
  entanglement entropy beyond the classical regime},  {\em JHEP 01 (2015) 073}
  (Aug., 2014) [\href{http://arxiv.org/abs/1408.3203}{{\tt arXiv:1408.3203}}].

\bibitem{Harper2018}
J.~Harper, M.~Headrick, and A.~Rolph, {\it Bit threads in higher curvature
  gravity},  {\em JHEP} {\bf 11} (2018) 168,
  [\href{http://arxiv.org/abs/1807.04294}{{\tt arXiv:1807.04294}}].

\bibitem{Rolph2021}
A.~Rolph, {\it Quantum bit threads},
  \href{http://arxiv.org/abs/2105.08072}{{\tt arXiv:2105.08072}}.

\bibitem{Almheiri2019b}
A.~Almheiri, R.~Mahajan, J.~Maldacena, and Y.~Zhao, {\it The page curve of
  hawking radiation from semiclassical geometry},
  \href{http://arxiv.org/abs/1908.10996}{{\tt arXiv:1908.10996}}.

\bibitem{Akers2019}
C.~Akers, N.~Engelhardt, G.~Penington, and M.~Usatyuk, {\it Quantum maximin
  surfaces},  \href{http://arxiv.org/abs/1912.02799}{{\tt arXiv:1912.02799}}.

\bibitem{Almheiri2021}
A.~Almheiri, T.~Hartman, J.~Maldacena, E.~Shaghoulian, and A.~Tajdini, {\it The
  entropy of hawking radiation},  {\em Rev. Mod. Phys.} {\bf 93} (2021), no.~3
  035002, [\href{http://arxiv.org/abs/2006.06872}{{\tt arXiv:2006.06872}}].

\bibitem{Almheiri2019}
A.~Almheiri, R.~Mahajan, and J.~Maldacena, {\it Islands outside the horizon},
  \href{http://arxiv.org/abs/1910.11077}{{\tt arXiv:1910.11077}}.

\bibitem{Almheiri2019a}
A.~Almheiri, N.~Engelhardt, D.~Marolf, and H.~Maxfield, {\it The entropy of
  bulk quantum fields and the entanglement wedge of an evaporating black hole},
   \href{http://arxiv.org/abs/1905.08762}{{\tt arXiv:1905.08762}}.

\bibitem{Almheiri2020}
A.~Almheiri, R.~Mahajan, and J.~E. Santos, {\it Entanglement islands in higher
  dimensions},  {\em SciPost Phys.} {\bf 9} (2020), no.~1 001,
  [\href{http://arxiv.org/abs/1911.09666}{{\tt arXiv:1911.09666}}].

\bibitem{Hartman2020}
T.~Hartman, Y.~Jiang, and E.~Shaghoulian, {\it Islands in cosmology},  {\em
  JHEP} {\bf 11} (2020) 111, [\href{http://arxiv.org/abs/2008.01022}{{\tt
  arXiv:2008.01022}}].

\bibitem{Bousso2015}
R.~Bousso, Z.~Fisher, S.~Leichenauer, , and A.~C. Wall, {\it A quantum
  focussing conjecture},  {\em Phys. Rev. D 93, 064044 (2016)} (June, 2015)
  [\href{http://arxiv.org/abs/1506.02669}{{\tt arXiv:1506.02669}}].

\bibitem{Casini2007}
H.~Casini, {\it Mutual information challenges entropy bounds},  {\em Class.
  Quant. Grav.} {\bf 24} (2007) 1293--1302,
  [\href{http://arxiv.org/abs/gr-qc/0609126}{{\tt gr-qc/0609126}}].

\bibitem{Bombelli1986}
L.~Bombelli, R.~K. Koul, J.~Lee, and R.~D. Sorkin, {\it A quantum source of
  entropy for black holes},  {\em Phys. Rev. D} {\bf 34} (1986) 373--383.

\bibitem{Wall2012}
A.~C. Wall, {\it Maximin surfaces, and the strong subadditivity of the
  covariant holographic entanglement entropy},  {\em Class. Quantum Grav. 31,
  225007 (2014)} (Nov., 2012) [\href{http://arxiv.org/abs/1211.3494}{{\tt
  arXiv:1211.3494}}].

\bibitem{Wald:1984rg}
R.~M. Wald, {\em {General Relativity}}.
\newblock Chicago Univ. Pr., Chicago, USA, 1984.

\bibitem{Casini2008}
H.~Casini, {\it Relative entropy and the bekenstein bound},  {\em Class. Quant.
  Grav.} {\bf 25} (2008) 205021, [\href{http://arxiv.org/abs/0804.2182}{{\tt
  arXiv:0804.2182}}].

\bibitem{Bousso2014}
R.~Bousso, H.~Casini, Z.~Fisher, and J.~Maldacena, {\it Proof of a quantum
  bousso bound},  {\em Phys. Rev. D} {\bf 90} (2014), no.~4 044002,
  [\href{http://arxiv.org/abs/1404.5635}{{\tt arXiv:1404.5635}}].

\bibitem{Strominger2003}
A.~Strominger and D.~Thompson, {\it A quantum bousso bound},  {\em
  Phys.Rev.D70:044007,2004} (Mar., 2003)
  [\href{http://arxiv.org/abs/hep-th/0303067}{{\tt hep-th/0303067}}].

\bibitem{BoussoTalk}
R.~Bousso, ``Quantum singularities.'' International Conference on Strings,
  Fields and Holograms at Monte Verità Conference Center, Ascona,
  Switzerland., Oct., 2021.

\bibitem{Raychaudhuri:1953yv}
A.~Raychaudhuri, {\it {Relativistic cosmology. 1.}},  {\em Phys. Rev.} {\bf 98}
  (1955) 1123--1126.

\bibitem{Lowe1999}
D.~A. Lowe, {\it Comments on a covariant entropy conjecture},  {\em JHEP} {\bf
  10} (1999) 026, [\href{http://arxiv.org/abs/hep-th/9907062}{{\tt
  hep-th/9907062}}].

\bibitem{Bousso1999}
R.~Bousso, {\it A covariant entropy conjecture},  {\em JHEP} {\bf 07} (1999)
  004, [\href{http://arxiv.org/abs/hep-th/9905177}{{\tt hep-th/9905177}}].

\bibitem{Akers2018}
C.~Akers, R.~Bousso, I.~F. Halpern, and G.~N. Remmen, {\it Boundary of the
  future of a surface},  {\em Phys. Rev. D} {\bf 97} (2018), no.~2 024018,
  [\href{http://arxiv.org/abs/1711.06689}{{\tt arXiv:1711.06689}}].

\bibitem{Jacobson2005}
T.~Jacobson, {\em Introduction to Quantum Fields in Curved Spacetime and the
  Hawking Effect}, pp.~39--89.
\newblock Springer US, Boston, MA, 2005.

\bibitem{Almheiri2013}
A.~Almheiri, D.~Marolf, J.~Polchinski, and J.~Sully, {\it Black holes:
  Complementarity or firewalls?},  {\em JHEP} {\bf 02} (2013) 062,
  [\href{http://arxiv.org/abs/1207.3123}{{\tt arXiv:1207.3123}}].

\bibitem{Bianchi2012}
E.~Bianchi and R.~C. Myers, {\it On the architecture of spacetime geometry},
  \href{http://arxiv.org/abs/1212.5183}{{\tt arXiv:1212.5183}}.

\bibitem{Cooperman2013}
J.~H. Cooperman and M.~A. Luty, {\it Renormalization of entanglement entropy
  and the gravitational effective action},  {\em Journal of High Energy Physics
  12 (2014) 045} (Feb., 2013) [\href{http://arxiv.org/abs/1302.1878}{{\tt
  arXiv:1302.1878}}].

\bibitem{Wall2009}
A.~C. Wall, {\it Ten proofs of the generalized second law},  {\em JHEP} {\bf
  06} (2009) 021, [\href{http://arxiv.org/abs/0901.3865}{{\tt
  arXiv:0901.3865}}].

\bibitem{Penington2019}
G.~Penington, {\it Entanglement wedge reconstruction and the information
  paradox},  \href{http://arxiv.org/abs/1905.08255}{{\tt arXiv:1905.08255}}.

\bibitem{Harlow2016}
D.~Harlow, {\it Jerusalem lectures on black holes and quantum information},
  {\em Rev. Mod. Phys.} {\bf 88} (2016) 015002,
  [\href{http://arxiv.org/abs/1409.1231}{{\tt arXiv:1409.1231}}].

\bibitem{Jacobson2003}
T.~Jacobson and R.~Parentani, {\it Horizon entropy},  {\em Found.Phys. 33
  (2003) 323-348} (Feb., 2003) [\href{http://arxiv.org/abs/gr-qc/0302099}{{\tt
  gr-qc/0302099}}].

\bibitem{Engelhardt2021}
N.~Engelhardt, G.~Penington, and A.~Shahbazi-Moghaddam, {\it Finding pythons in
  unexpected places},  \href{http://arxiv.org/abs/2105.09316}{{\tt
  arXiv:2105.09316}}.

\bibitem{Flanagan2000}
E.~E. Flanagan, D.~Marolf, and R.~M. Wald, {\it Proof of classical versions of
  the bousso entropy bound and of the generalized second law},  {\em Phys. Rev.
  D} {\bf 62} (2000) 084035, [\href{http://arxiv.org/abs/hep-th/9908070}{{\tt
  hep-th/9908070}}].

\bibitem{Bousso2000}
R.~Bousso, {\it The holographic principle for general backgrounds},  {\em
  Class. Quant. Grav.} {\bf 17} (2000) 997--1005,
  [\href{http://arxiv.org/abs/hep-th/9911002}{{\tt hep-th/9911002}}].

\bibitem{Bousso2002}
R.~Bousso, {\it The holographic principle},  {\em Rev. Mod. Phys.} {\bf 74}
  (2002) 825--874, [\href{http://arxiv.org/abs/hep-th/0203101}{{\tt
  hep-th/0203101}}].

\end{thebibliography}\endgroup
